% paper57-semantic-search.tex — Using the Semantic Site as an Index for Code Search and Retrieval
%   Paper 57 of the Judgment Geometry series.
% Compile: pdflatex paper57-semantic-search
\documentclass[11pt]{article}
\usepackage{lmodern}
% jugeo-common.tex — shared preamble for all Judgment Geometry papers
% Include via: % jugeo-common.tex — shared preamble for all Judgment Geometry papers
% Include via: % jugeo-common.tex — shared preamble for all Judgment Geometry papers
% Include via: \input{jugeo-common}

% ── Packages ────────────────────────────────────────────────────────────
\usepackage{amsmath,amssymb,amsthm,mathtools}
\usepackage{tikz-cd}
\usepackage{listings}
\usepackage{xcolor}
\usepackage{xspace}
\usepackage{booktabs}
\usepackage{multirow}
\usepackage{enumitem}
\usepackage[expansion=false]{microtype}
\usepackage{hyperref}
\usepackage{cleveref}
% \usepackage{thmtools}  % disabled: not available in this TeX installation
\usepackage{float}
\usepackage{caption}
\usepackage{subcaption}
\usepackage{tabularx}
\usepackage{stmaryrd}       % \llbracket, \rrbracket
\usepackage{bbm}            % \mathbbm{1}

% ── Page geometry ───────────────────────────────────────────────────────
\usepackage[margin=1in]{geometry}

% ── Theorem environments ────────────────────────────────────────────────
\theoremstyle{plain}
\newtheorem{theorem}{Theorem}[section]
\newtheorem{lemma}[theorem]{Lemma}
\newtheorem{proposition}[theorem]{Proposition}
\newtheorem{corollary}[theorem]{Corollary}

\theoremstyle{definition}
\newtheorem{definition}[theorem]{Definition}
\newtheorem{example}[theorem]{Example}
\newtheorem{construction}[theorem]{Construction}

\theoremstyle{remark}
\newtheorem{remark}[theorem]{Remark}
\newtheorem{notation}[theorem]{Notation}

% ── Python listings style ──────────────────────────────────────────────
\definecolor{jugeoBlue}{HTML}{2563EB}
\definecolor{jugeoGreen}{HTML}{059669}
\definecolor{jugeoRed}{HTML}{DC2626}
\definecolor{jugeoPurple}{HTML}{7C3AED}
\definecolor{jugeoGray}{HTML}{6B7280}
\definecolor{jugeoBg}{HTML}{F8FAFC}

\lstdefinestyle{jugeo-python}{
  language=Python,
  basicstyle=\ttfamily\small,
  keywordstyle=\color{jugeoBlue}\bfseries,
  stringstyle=\color{jugeoGreen},
  commentstyle=\color{jugeoGray}\itshape,
  numberstyle=\tiny\color{jugeoGray},
  backgroundcolor=\color{jugeoBg},
  numbers=left,
  numbersep=6pt,
  frame=single,
  framerule=0.4pt,
  rulecolor=\color{jugeoGray!40},
  breaklines=true,
  breakatwhitespace=true,
  showstringspaces=false,
  tabsize=4,
  xleftmargin=1.5em,
  framexleftmargin=1.5em,
  morekeywords={dataclass,frozen,field,Enum,IntEnum,auto,Any,
                Mapping,Sequence,Optional,match,case},
  literate={→}{{$\to$}}1 {←}{{$\leftarrow$}}1
           {≤}{{$\leq$}}1 {≥}{{$\geq$}}1 {≠}{{$\neq$}}1
           {∈}{{$\in$}}1 {∀}{{$\forall$}}1 {∃}{{$\exists$}}1
           {⊕}{{$\oplus$}}1 {⊖}{{$\ominus$}}1,
  escapeinside={(*@}{@*)},
}
\lstset{style=jugeo-python}

\lstdefinestyle{jugeo-output}{
  basicstyle=\ttfamily\small,
  backgroundcolor=\color{jugeoBg},
  frame=single,
  framerule=0.4pt,
  rulecolor=\color{jugeoGray!40},
  breaklines=true,
  showstringspaces=false,
  xleftmargin=1.5em,
  framexleftmargin=0.5em,
  numbers=none,
}

% ── JuGeo-specific macros ──────────────────────────────────────────────

% Judgment 8-tuple
\newcommand{\J}{\mathcal{J}}
\newcommand{\Jud}[8]{(#1,\,#2,\,#3,\,#4,\,#5,\,#6,\,#7,\,#8)}
\newcommand{\JudShort}[1]{\J_{#1}}

% Components
\newcommand{\coord}{\mathsf{c}}            % coordinate
\newcommand{\prop}{\varphi}                 % proposition
\newcommand{\carrier}{\mathcal{A}}          % carrier type
\newcommand{\evid}{\mathcal{E}}             % evidence bundle
\newcommand{\oblig}{\mathcal{O}}            % obligations
\newcommand{\obstr}{\mathcal{B}}            % obstructions
\newcommand{\trust}{\mathcal{T}}            % trust annotation
\newcommand{\prov}{\Pi}                     % provenance

% Trust algebra
\newcommand{\Talg}{\mathfrak{T}}            % trust ordered algebra
\newcommand{\Eadm}{\mathcal{E}_{\mathrm{adm}}}
\newcommand{\tleq}{\preceq}                 % trust ordering
\newcommand{\tjoin}{\oplus}                 % conservative join
\newcommand{\tatten}{\ominus}               % attenuation
\newcommand{\tpromote}{\uparrow_\pi}        % promotion
\newcommand{\tdemote}{\downarrow_\chi}       % demotion

% Trust tiers
\newcommand{\Tcontra}{\ensuremath{\mathsf{CONTRADICTED}}}
\newcommand{\Tunver}{\ensuremath{\mathsf{UNVERIFIED}}}
\newcommand{\Tcopilot}{\ensuremath{\mathsf{COPILOT}}}
\newcommand{\Toracle}{\ensuremath{\mathsf{ORACLE}}}
\newcommand{\Truntime}{\ensuremath{\mathsf{RUNTIME}}}
\newcommand{\Tsolver}{\ensuremath{\mathsf{SOLVER}}}
\newcommand{\Tproof}{\ensuremath{\mathsf{PROOF}}}

% Site and geometry
\newcommand{\Site}{\mathcal{S}}             % semantic site
\newcommand{\Cov}{\mathrm{Cov}}            % covering family
\newcommand{\Ob}{\mathrm{Ob}}              % objects
\newcommand{\Mor}{\mathrm{Mor}}            % morphisms
\newcommand{\res}{\mathrm{res}}            % restriction
\newcommand{\Cech}{\check{C}}              % Čech complex
\newcommand{\Hcoh}[1]{H^{#1}}             % cohomology group

% Descent
\newcommand{\descent}{\mathrm{desc}}
\newcommand{\glue}{\mathrm{glue}}
\newcommand{\overlap}[2]{#1 \cap #2}

% Semantic moves
\newcommand{\VERIFY}{\textsc{Verify}}
\newcommand{\CONSTRUCT}{\textsc{Construct}}
\newcommand{\NORMALIZE}{\textsc{Normalize}}
\newcommand{\GLUE}{\textsc{Glue}}
\newcommand{\RESTRICT}{\textsc{Restrict}}
\newcommand{\TRANSPORT}{\textsc{Transport}}
\newcommand{\REPAIR}{\textsc{Repair}}
\newcommand{\PROMOTE}{\textsc{Promote}}
\newcommand{\CHALLENGE}{\textsc{Challenge}}
\newcommand{\ESCALATE}{\textsc{Escalate}}

% Fragments
\newcommand{\QFLIA}{\texttt{QF\_LIA}}
\newcommand{\QFLRA}{\texttt{QF\_LRA}}
\newcommand{\QFBV}{\texttt{QF\_BV}}
\newcommand{\QFUF}{\texttt{QF\_UF}}

% Misc
\newcommand{\Python}{\textsc{Python}}
\newcommand{\JuGeo}{\textsc{JuGeo}}
\newcommand{\LEAN}{\textsc{Lean}\,4}
\newcommand{\Fstar}{F$^{\star}$}
\newcommand{\Coq}{\textsc{Coq}}
\newcommand{\Dafny}{\textsc{Dafny}}

% Common shortcuts
\newcommand{\ie}{i.e.\@\xspace}
\newcommand{\eg}{e.g.\@\xspace}
\newcommand{\cf}{cf.\@\xspace}
\newcommand{\wrt}{w.r.t.\@\xspace}

% Evaluation shortcuts
\newcommand{\numcases}{300}
\newcommand{\numspec}{100}
\newcommand{\numeq}{100}
\newcommand{\numbug}{100}
\newcommand{\medtime}{6.5\,ms}
\newcommand{\totaltime}{1.949\,s}


% ── Packages ────────────────────────────────────────────────────────────
\usepackage{amsmath,amssymb,amsthm,mathtools}
\usepackage{tikz-cd}
\usepackage{listings}
\usepackage{xcolor}
\usepackage{xspace}
\usepackage{booktabs}
\usepackage{multirow}
\usepackage{enumitem}
\usepackage[expansion=false]{microtype}
\usepackage{hyperref}
\usepackage{cleveref}
% \usepackage{thmtools}  % disabled: not available in this TeX installation
\usepackage{float}
\usepackage{caption}
\usepackage{subcaption}
\usepackage{tabularx}
\usepackage{stmaryrd}       % \llbracket, \rrbracket
\usepackage{bbm}            % \mathbbm{1}

% ── Page geometry ───────────────────────────────────────────────────────
\usepackage[margin=1in]{geometry}

% ── Theorem environments ────────────────────────────────────────────────
\theoremstyle{plain}
\newtheorem{theorem}{Theorem}[section]
\newtheorem{lemma}[theorem]{Lemma}
\newtheorem{proposition}[theorem]{Proposition}
\newtheorem{corollary}[theorem]{Corollary}

\theoremstyle{definition}
\newtheorem{definition}[theorem]{Definition}
\newtheorem{example}[theorem]{Example}
\newtheorem{construction}[theorem]{Construction}

\theoremstyle{remark}
\newtheorem{remark}[theorem]{Remark}
\newtheorem{notation}[theorem]{Notation}

% ── Python listings style ──────────────────────────────────────────────
\definecolor{jugeoBlue}{HTML}{2563EB}
\definecolor{jugeoGreen}{HTML}{059669}
\definecolor{jugeoRed}{HTML}{DC2626}
\definecolor{jugeoPurple}{HTML}{7C3AED}
\definecolor{jugeoGray}{HTML}{6B7280}
\definecolor{jugeoBg}{HTML}{F8FAFC}

\lstdefinestyle{jugeo-python}{
  language=Python,
  basicstyle=\ttfamily\small,
  keywordstyle=\color{jugeoBlue}\bfseries,
  stringstyle=\color{jugeoGreen},
  commentstyle=\color{jugeoGray}\itshape,
  numberstyle=\tiny\color{jugeoGray},
  backgroundcolor=\color{jugeoBg},
  numbers=left,
  numbersep=6pt,
  frame=single,
  framerule=0.4pt,
  rulecolor=\color{jugeoGray!40},
  breaklines=true,
  breakatwhitespace=true,
  showstringspaces=false,
  tabsize=4,
  xleftmargin=1.5em,
  framexleftmargin=1.5em,
  morekeywords={dataclass,frozen,field,Enum,IntEnum,auto,Any,
                Mapping,Sequence,Optional,match,case},
  literate={→}{{$\to$}}1 {←}{{$\leftarrow$}}1
           {≤}{{$\leq$}}1 {≥}{{$\geq$}}1 {≠}{{$\neq$}}1
           {∈}{{$\in$}}1 {∀}{{$\forall$}}1 {∃}{{$\exists$}}1
           {⊕}{{$\oplus$}}1 {⊖}{{$\ominus$}}1,
  escapeinside={(*@}{@*)},
}
\lstset{style=jugeo-python}

\lstdefinestyle{jugeo-output}{
  basicstyle=\ttfamily\small,
  backgroundcolor=\color{jugeoBg},
  frame=single,
  framerule=0.4pt,
  rulecolor=\color{jugeoGray!40},
  breaklines=true,
  showstringspaces=false,
  xleftmargin=1.5em,
  framexleftmargin=0.5em,
  numbers=none,
}

% ── JuGeo-specific macros ──────────────────────────────────────────────

% Judgment 8-tuple
\newcommand{\J}{\mathcal{J}}
\newcommand{\Jud}[8]{(#1,\,#2,\,#3,\,#4,\,#5,\,#6,\,#7,\,#8)}
\newcommand{\JudShort}[1]{\J_{#1}}

% Components
\newcommand{\coord}{\mathsf{c}}            % coordinate
\newcommand{\prop}{\varphi}                 % proposition
\newcommand{\carrier}{\mathcal{A}}          % carrier type
\newcommand{\evid}{\mathcal{E}}             % evidence bundle
\newcommand{\oblig}{\mathcal{O}}            % obligations
\newcommand{\obstr}{\mathcal{B}}            % obstructions
\newcommand{\trust}{\mathcal{T}}            % trust annotation
\newcommand{\prov}{\Pi}                     % provenance

% Trust algebra
\newcommand{\Talg}{\mathfrak{T}}            % trust ordered algebra
\newcommand{\Eadm}{\mathcal{E}_{\mathrm{adm}}}
\newcommand{\tleq}{\preceq}                 % trust ordering
\newcommand{\tjoin}{\oplus}                 % conservative join
\newcommand{\tatten}{\ominus}               % attenuation
\newcommand{\tpromote}{\uparrow_\pi}        % promotion
\newcommand{\tdemote}{\downarrow_\chi}       % demotion

% Trust tiers
\newcommand{\Tcontra}{\ensuremath{\mathsf{CONTRADICTED}}}
\newcommand{\Tunver}{\ensuremath{\mathsf{UNVERIFIED}}}
\newcommand{\Tcopilot}{\ensuremath{\mathsf{COPILOT}}}
\newcommand{\Toracle}{\ensuremath{\mathsf{ORACLE}}}
\newcommand{\Truntime}{\ensuremath{\mathsf{RUNTIME}}}
\newcommand{\Tsolver}{\ensuremath{\mathsf{SOLVER}}}
\newcommand{\Tproof}{\ensuremath{\mathsf{PROOF}}}

% Site and geometry
\newcommand{\Site}{\mathcal{S}}             % semantic site
\newcommand{\Cov}{\mathrm{Cov}}            % covering family
\newcommand{\Ob}{\mathrm{Ob}}              % objects
\newcommand{\Mor}{\mathrm{Mor}}            % morphisms
\newcommand{\res}{\mathrm{res}}            % restriction
\newcommand{\Cech}{\check{C}}              % Čech complex
\newcommand{\Hcoh}[1]{H^{#1}}             % cohomology group

% Descent
\newcommand{\descent}{\mathrm{desc}}
\newcommand{\glue}{\mathrm{glue}}
\newcommand{\overlap}[2]{#1 \cap #2}

% Semantic moves
\newcommand{\VERIFY}{\textsc{Verify}}
\newcommand{\CONSTRUCT}{\textsc{Construct}}
\newcommand{\NORMALIZE}{\textsc{Normalize}}
\newcommand{\GLUE}{\textsc{Glue}}
\newcommand{\RESTRICT}{\textsc{Restrict}}
\newcommand{\TRANSPORT}{\textsc{Transport}}
\newcommand{\REPAIR}{\textsc{Repair}}
\newcommand{\PROMOTE}{\textsc{Promote}}
\newcommand{\CHALLENGE}{\textsc{Challenge}}
\newcommand{\ESCALATE}{\textsc{Escalate}}

% Fragments
\newcommand{\QFLIA}{\texttt{QF\_LIA}}
\newcommand{\QFLRA}{\texttt{QF\_LRA}}
\newcommand{\QFBV}{\texttt{QF\_BV}}
\newcommand{\QFUF}{\texttt{QF\_UF}}

% Misc
\newcommand{\Python}{\textsc{Python}}
\newcommand{\JuGeo}{\textsc{JuGeo}}
\newcommand{\LEAN}{\textsc{Lean}\,4}
\newcommand{\Fstar}{F$^{\star}$}
\newcommand{\Coq}{\textsc{Coq}}
\newcommand{\Dafny}{\textsc{Dafny}}

% Common shortcuts
\newcommand{\ie}{i.e.\@\xspace}
\newcommand{\eg}{e.g.\@\xspace}
\newcommand{\cf}{cf.\@\xspace}
\newcommand{\wrt}{w.r.t.\@\xspace}

% Evaluation shortcuts
\newcommand{\numcases}{300}
\newcommand{\numspec}{100}
\newcommand{\numeq}{100}
\newcommand{\numbug}{100}
\newcommand{\medtime}{6.5\,ms}
\newcommand{\totaltime}{1.949\,s}


% ── Packages ────────────────────────────────────────────────────────────
\usepackage{amsmath,amssymb,amsthm,mathtools}
\usepackage{tikz-cd}
\usepackage{listings}
\usepackage{xcolor}
\usepackage{xspace}
\usepackage{booktabs}
\usepackage{multirow}
\usepackage{enumitem}
\usepackage[expansion=false]{microtype}
\usepackage{hyperref}
\usepackage{cleveref}
% \usepackage{thmtools}  % disabled: not available in this TeX installation
\usepackage{float}
\usepackage{caption}
\usepackage{subcaption}
\usepackage{tabularx}
\usepackage{stmaryrd}       % \llbracket, \rrbracket
\usepackage{bbm}            % \mathbbm{1}

% ── Page geometry ───────────────────────────────────────────────────────
\usepackage[margin=1in]{geometry}

% ── Theorem environments ────────────────────────────────────────────────
\theoremstyle{plain}
\newtheorem{theorem}{Theorem}[section]
\newtheorem{lemma}[theorem]{Lemma}
\newtheorem{proposition}[theorem]{Proposition}
\newtheorem{corollary}[theorem]{Corollary}

\theoremstyle{definition}
\newtheorem{definition}[theorem]{Definition}
\newtheorem{example}[theorem]{Example}
\newtheorem{construction}[theorem]{Construction}

\theoremstyle{remark}
\newtheorem{remark}[theorem]{Remark}
\newtheorem{notation}[theorem]{Notation}

% ── Python listings style ──────────────────────────────────────────────
\definecolor{jugeoBlue}{HTML}{2563EB}
\definecolor{jugeoGreen}{HTML}{059669}
\definecolor{jugeoRed}{HTML}{DC2626}
\definecolor{jugeoPurple}{HTML}{7C3AED}
\definecolor{jugeoGray}{HTML}{6B7280}
\definecolor{jugeoBg}{HTML}{F8FAFC}

\lstdefinestyle{jugeo-python}{
  language=Python,
  basicstyle=\ttfamily\small,
  keywordstyle=\color{jugeoBlue}\bfseries,
  stringstyle=\color{jugeoGreen},
  commentstyle=\color{jugeoGray}\itshape,
  numberstyle=\tiny\color{jugeoGray},
  backgroundcolor=\color{jugeoBg},
  numbers=left,
  numbersep=6pt,
  frame=single,
  framerule=0.4pt,
  rulecolor=\color{jugeoGray!40},
  breaklines=true,
  breakatwhitespace=true,
  showstringspaces=false,
  tabsize=4,
  xleftmargin=1.5em,
  framexleftmargin=1.5em,
  morekeywords={dataclass,frozen,field,Enum,IntEnum,auto,Any,
                Mapping,Sequence,Optional,match,case},
  literate={→}{{$\to$}}1 {←}{{$\leftarrow$}}1
           {≤}{{$\leq$}}1 {≥}{{$\geq$}}1 {≠}{{$\neq$}}1
           {∈}{{$\in$}}1 {∀}{{$\forall$}}1 {∃}{{$\exists$}}1
           {⊕}{{$\oplus$}}1 {⊖}{{$\ominus$}}1,
  escapeinside={(*@}{@*)},
}
\lstset{style=jugeo-python}

\lstdefinestyle{jugeo-output}{
  basicstyle=\ttfamily\small,
  backgroundcolor=\color{jugeoBg},
  frame=single,
  framerule=0.4pt,
  rulecolor=\color{jugeoGray!40},
  breaklines=true,
  showstringspaces=false,
  xleftmargin=1.5em,
  framexleftmargin=0.5em,
  numbers=none,
}

% ── JuGeo-specific macros ──────────────────────────────────────────────

% Judgment 8-tuple
\newcommand{\J}{\mathcal{J}}
\newcommand{\Jud}[8]{(#1,\,#2,\,#3,\,#4,\,#5,\,#6,\,#7,\,#8)}
\newcommand{\JudShort}[1]{\J_{#1}}

% Components
\newcommand{\coord}{\mathsf{c}}            % coordinate
\newcommand{\prop}{\varphi}                 % proposition
\newcommand{\carrier}{\mathcal{A}}          % carrier type
\newcommand{\evid}{\mathcal{E}}             % evidence bundle
\newcommand{\oblig}{\mathcal{O}}            % obligations
\newcommand{\obstr}{\mathcal{B}}            % obstructions
\newcommand{\trust}{\mathcal{T}}            % trust annotation
\newcommand{\prov}{\Pi}                     % provenance

% Trust algebra
\newcommand{\Talg}{\mathfrak{T}}            % trust ordered algebra
\newcommand{\Eadm}{\mathcal{E}_{\mathrm{adm}}}
\newcommand{\tleq}{\preceq}                 % trust ordering
\newcommand{\tjoin}{\oplus}                 % conservative join
\newcommand{\tatten}{\ominus}               % attenuation
\newcommand{\tpromote}{\uparrow_\pi}        % promotion
\newcommand{\tdemote}{\downarrow_\chi}       % demotion

% Trust tiers
\newcommand{\Tcontra}{\ensuremath{\mathsf{CONTRADICTED}}}
\newcommand{\Tunver}{\ensuremath{\mathsf{UNVERIFIED}}}
\newcommand{\Tcopilot}{\ensuremath{\mathsf{COPILOT}}}
\newcommand{\Toracle}{\ensuremath{\mathsf{ORACLE}}}
\newcommand{\Truntime}{\ensuremath{\mathsf{RUNTIME}}}
\newcommand{\Tsolver}{\ensuremath{\mathsf{SOLVER}}}
\newcommand{\Tproof}{\ensuremath{\mathsf{PROOF}}}

% Site and geometry
\newcommand{\Site}{\mathcal{S}}             % semantic site
\newcommand{\Cov}{\mathrm{Cov}}            % covering family
\newcommand{\Ob}{\mathrm{Ob}}              % objects
\newcommand{\Mor}{\mathrm{Mor}}            % morphisms
\newcommand{\res}{\mathrm{res}}            % restriction
\newcommand{\Cech}{\check{C}}              % Čech complex
\newcommand{\Hcoh}[1]{H^{#1}}             % cohomology group

% Descent
\newcommand{\descent}{\mathrm{desc}}
\newcommand{\glue}{\mathrm{glue}}
\newcommand{\overlap}[2]{#1 \cap #2}

% Semantic moves
\newcommand{\VERIFY}{\textsc{Verify}}
\newcommand{\CONSTRUCT}{\textsc{Construct}}
\newcommand{\NORMALIZE}{\textsc{Normalize}}
\newcommand{\GLUE}{\textsc{Glue}}
\newcommand{\RESTRICT}{\textsc{Restrict}}
\newcommand{\TRANSPORT}{\textsc{Transport}}
\newcommand{\REPAIR}{\textsc{Repair}}
\newcommand{\PROMOTE}{\textsc{Promote}}
\newcommand{\CHALLENGE}{\textsc{Challenge}}
\newcommand{\ESCALATE}{\textsc{Escalate}}

% Fragments
\newcommand{\QFLIA}{\texttt{QF\_LIA}}
\newcommand{\QFLRA}{\texttt{QF\_LRA}}
\newcommand{\QFBV}{\texttt{QF\_BV}}
\newcommand{\QFUF}{\texttt{QF\_UF}}

% Misc
\newcommand{\Python}{\textsc{Python}}
\newcommand{\JuGeo}{\textsc{JuGeo}}
\newcommand{\LEAN}{\textsc{Lean}\,4}
\newcommand{\Fstar}{F$^{\star}$}
\newcommand{\Coq}{\textsc{Coq}}
\newcommand{\Dafny}{\textsc{Dafny}}

% Common shortcuts
\newcommand{\ie}{i.e.\@\xspace}
\newcommand{\eg}{e.g.\@\xspace}
\newcommand{\cf}{cf.\@\xspace}
\newcommand{\wrt}{w.r.t.\@\xspace}

% Evaluation shortcuts
\newcommand{\numcases}{300}
\newcommand{\numspec}{100}
\newcommand{\numeq}{100}
\newcommand{\numbug}{100}
\newcommand{\medtime}{6.5\,ms}
\newcommand{\totaltime}{1.949\,s}

% experiment-data.tex — AUTO-GENERATED from real jugeo CLI runs
% DO NOT EDIT — regenerate with: python3 experiments/run_universal_metrics.py
% Generated from 30 programs across 6 domains

\newcommand{\expTotalPrograms}{30}
\newcommand{\expVerified}{30}
\newcommand{\expAccuracy}{100.0\%}
\newcommand{\expCoordsMin}{2}
\newcommand{\expCoordsMax}{10}
\newcommand{\expCoordsMean}{5.2}
\newcommand{\expCoordsMedian}{5.5}
\newcommand{\expPropsMin}{4}
\newcommand{\expPropsMax}{20}
\newcommand{\expPropsMean}{10.4}
\newcommand{\expPropsSum}{311}
\newcommand{\expPropsOkSum}{310}
\newcommand{\expTimeMin}{3.506\,s}
\newcommand{\expTimeMax}{6.714\,s}
\newcommand{\expTimeMean}{4.64\,s}
\newcommand{\expTimeMedian}{4.508\,s}
\newcommand{\expTimeTotal}{139.204\,s}
\newcommand{\expObstructionTotal}{0}
\newcommand{\expHOneZero}{30}
\newcommand{\expDomainCount}{6}
\newcommand{\expTrustCOPILOTSUGGESTED}{30}
\newcommand{\expDomainDsCount}{5}
\newcommand{\expDomainDsVerified}{5}
\newcommand{\expDomainDsAvgCoords}{7.6}
\newcommand{\expDomainDsAvgProps}{15.2}
\newcommand{\expDomainMathCount}{5}
\newcommand{\expDomainMathVerified}{5}
\newcommand{\expDomainMathAvgCoords}{4.8}
\newcommand{\expDomainMathAvgProps}{9.4}
\newcommand{\expDomainSmCount}{5}
\newcommand{\expDomainSmVerified}{5}
\newcommand{\expDomainSmAvgCoords}{7.6}
\newcommand{\expDomainSmAvgProps}{15.2}
\newcommand{\expDomainSortCount}{5}
\newcommand{\expDomainSortVerified}{5}
\newcommand{\expDomainSortAvgCoords}{2.4}
\newcommand{\expDomainSortAvgProps}{4.8}
\newcommand{\expDomainStrCount}{5}
\newcommand{\expDomainStrVerified}{5}
\newcommand{\expDomainStrAvgCoords}{3.2}
\newcommand{\expDomainStrAvgProps}{6.4}
\newcommand{\expDomainWebCount}{5}
\newcommand{\expDomainWebVerified}{5}
\newcommand{\expDomainWebAvgCoords}{5.6}
\newcommand{\expDomainWebAvgProps}{11.2}

% subsystem-data.tex — AUTO-GENERATED from real jugeo subsystem experiments
% DO NOT EDIT — regenerate with: python3 experiments/run_subsystem_experiments.py

\newcommand{\subCliTotal}{10}
\newcommand{\subCliVerified}{9}
\newcommand{\subCliAccuracy}{90.0\%}
\newcommand{\subCliMeanTime}{4.132\,s}
\newcommand{\subCliMinTime}{3.472\,s}
\newcommand{\subCliMaxTime}{5.372\,s}
\newcommand{\subCliTotalCoords}{54}
\newcommand{\subCliTotalProps}{107}
\newcommand{\subCliTotalPropsOk}{106}
\newcommand{\subSiteCalls}{90}
\newcommand{\subSiteSuccessful}{69}
\newcommand{\subSiteSuccessRate}{76.7\%}
\newcommand{\subSiteMeanTime}{0.0001\,s}
\newcommand{\subSiteTotalTime}{0.005\,s}
\newcommand{\subSpecificationsatisfactionSmallTime}{0.0003\,s}
\newcommand{\subSpecificationsatisfactionMediumTime}{0.0001\,s}
\newcommand{\subSpecificationsatisfactionLargeTime}{0.0001\,s}
\newcommand{\subInhabitantfleetSmallTime}{0.0\,s}
\newcommand{\subInhabitantfleetMediumTime}{0.0\,s}
\newcommand{\subInhabitantfleetLargeTime}{0.0\,s}
\newcommand{\subTheoremecologySmallTime}{0.0\,s}
\newcommand{\subTheoremecologyMediumTime}{0.0\,s}
\newcommand{\subTheoremecologyLargeTime}{0.0\,s}
\newcommand{\subStatespaceexplorationSmallTime}{0.0\,s}
\newcommand{\subStatespaceexplorationMediumTime}{0.0\,s}
\newcommand{\subStatespaceexplorationLargeTime}{0.0\,s}
\newcommand{\subRepairsemanticsSmallTime}{0.0\,s}
\newcommand{\subRepairsemanticsMediumTime}{0.0\,s}
\newcommand{\subRepairsemanticsLargeTime}{0.0\,s}
\newcommand{\subReplaygluingSmallTime}{0.0\,s}
\newcommand{\subReplaygluingMediumTime}{0.0\,s}
\newcommand{\subReplaygluingLargeTime}{0.0\,s}
\newcommand{\subSemanticfuturesSmallTime}{0.0\,s}
\newcommand{\subSemanticfuturesMediumTime}{0.0\,s}
\newcommand{\subSemanticfuturesLargeTime}{0.0\,s}
\newcommand{\subHypercovertreatySmallTime}{0.0\,s}
\newcommand{\subHypercovertreatyMediumTime}{0.0\,s}
\newcommand{\subHypercovertreatyLargeTime}{0.0\,s}
\newcommand{\subEncodeforsolverSmallTime}{0.0026\,s}
\newcommand{\subEncodeforsolverMediumTime}{0.0002\,s}
\newcommand{\subEncodeforsolverLargeTime}{0.0001\,s}
\newcommand{\subInterfaceroutingSmallTime}{0.0\,s}
\newcommand{\subInterfaceroutingMediumTime}{0.0\,s}
\newcommand{\subInterfaceroutingLargeTime}{0.0\,s}
\newcommand{\subOrchestrateverificationSmallTime}{0.0002\,s}
\newcommand{\subOrchestrateverificationMediumTime}{0.0\,s}
\newcommand{\subOrchestrateverificationLargeTime}{0.0\,s}
\newcommand{\subRunfulldescentSmallTime}{0.0\,s}
\newcommand{\subRunfulldescentMediumTime}{0.0\,s}
\newcommand{\subRunfulldescentLargeTime}{0.0\,s}
\newcommand{\subKernellifecycleSmallTime}{0.0\,s}
\newcommand{\subKernellifecycleMediumTime}{0.0\,s}
\newcommand{\subKernellifecycleLargeTime}{0.0\,s}
\newcommand{\subDiscoverypipelineSmallTime}{0.0\,s}
\newcommand{\subDiscoverypipelineMediumTime}{0.0\,s}
\newcommand{\subDiscoverypipelineLargeTime}{0.0\,s}
\newcommand{\subAnalogytransportSmallTime}{0.0\,s}
\newcommand{\subAnalogytransportMediumTime}{0.0\,s}
\newcommand{\subAnalogytransportLargeTime}{0.0\,s}
\newcommand{\subFormalcoresiteSmallTime}{0.0001\,s}
\newcommand{\subFormalcoresiteMediumTime}{0.0\,s}
\newcommand{\subFormalcoresiteLargeTime}{0.0\,s}
\newcommand{\subProblematlasSmallTime}{0.0\,s}
\newcommand{\subProblematlasMediumTime}{0.0\,s}
\newcommand{\subProblematlasLargeTime}{0.0\,s}
\newcommand{\subPublicalignmentSmallTime}{0.0\,s}
\newcommand{\subPublicalignmentMediumTime}{0.0\,s}
\newcommand{\subPublicalignmentLargeTime}{0.0\,s}
\newcommand{\subRegimebootstrappingSmallTime}{0.0\,s}
\newcommand{\subRegimebootstrappingMediumTime}{0.0\,s}
\newcommand{\subRegimebootstrappingLargeTime}{0.0\,s}
\newcommand{\subRelationalrefinementSmallTime}{0.0\,s}
\newcommand{\subRelationalrefinementMediumTime}{0.0\,s}
\newcommand{\subRelationalrefinementLargeTime}{0.0\,s}
\newcommand{\subSemanticclosureSmallTime}{0.0\,s}
\newcommand{\subSemanticclosureMediumTime}{0.0\,s}
\newcommand{\subSemanticclosureLargeTime}{0.0\,s}
\newcommand{\subTheoremeconomicsSmallTime}{0.0001\,s}
\newcommand{\subTheoremeconomicsMediumTime}{0.0\,s}
\newcommand{\subTheoremeconomicsLargeTime}{0.0\,s}
\newcommand{\subBenchmarksuiteSmallTime}{0.0\,s}
\newcommand{\subBenchmarksuiteMediumTime}{0.0\,s}
\newcommand{\subBenchmarksuiteLargeTime}{0.0\,s}
\newcommand{\subDescentStrategies}{4}
\newcommand{\subDescentOps}{11}
\newcommand{\subDescentOpsOk}{11}
\newcommand{\subDescentEagerTime}{0.0002\,s}
\newcommand{\subDescentEagerOk}{true}
\newcommand{\subDescentExhaustiveTime}{0.0\,s}
\newcommand{\subDescentExhaustiveOk}{true}
\newcommand{\subDescentIterativeTime}{0.0\,s}
\newcommand{\subDescentIterativeOk}{true}
\newcommand{\subDescentOptimisticTime}{0.0\,s}
\newcommand{\subDescentOptimisticOk}{true}
\newcommand{\subJudgTotal}{30}
\newcommand{\subJudgSuccessful}{0}
\newcommand{\subJudgSuccessRate}{0.0\%}
\newcommand{\subJudgMeanTimeUs}{7.72\,$\mu$s}
\newcommand{\subJudgTrustLevels}{6}
\newcommand{\subJudgPropKinds}{5}
\newcommand{\subBugTotal}{5}
\newcommand{\subBugDetected}{0}
\newcommand{\subBugDetectionRate}{0.0\%}
\newcommand{\subBugFalsePositiveRate}{0.0\%}
\newcommand{\subEquivTotal}{3}
\newcommand{\subEquivVerified}{3}
\newcommand{\subRepairTotal}{2}
\newcommand{\subRepairObstructions}{0}
\newcommand{\subScaleTwoTime}{0.0001\,s}
\newcommand{\subScaleFourTime}{0.0\,s}
\newcommand{\subScaleEightTime}{0.0\,s}
\newcommand{\subScaleTwelveTime}{0.0\,s}
\newcommand{\subScaleSixteenTime}{0.0\,s}
\newcommand{\subScaleTwentyTime}{0.0\,s}
\newcommand{\subTotalExperimentTime}{95.6\,s}

% comprehensive-data.tex — AUTO-GENERATED from real jugeo experiments
% DO NOT EDIT — regenerate with: python3 experiments/run_comprehensive_experiments.py
% Generated: 2026-03-23 09:19:06

% --- Size scaling metrics ---
\newcommand{\compTotalPrograms}{18}
\newcommand{\compTotalVerified}{18}
\newcommand{\compOverallAccuracy}{100.0\%}
\newcommand{\compTinyCount}{5}
\newcommand{\compTinyVerified}{5}
\newcommand{\compTinyMeanCoords}{3}
\newcommand{\compTinyMeanProps}{6}
\newcommand{\compTinyMeanTime}{4.95\,s}
\newcommand{\compTinyMeanLines}{5}
\newcommand{\compTinyMinTime}{4.45\,s}
\newcommand{\compTinyMaxTime}{5.51\,s}
\newcommand{\compSmallCount}{5}
\newcommand{\compSmallVerified}{5}
\newcommand{\compSmallMeanCoords}{3.6}
\newcommand{\compSmallMeanProps}{7.2}
\newcommand{\compSmallMeanTime}{4.85\,s}
\newcommand{\compSmallMeanLines}{16.0}
\newcommand{\compSmallMinTime}{4.30\,s}
\newcommand{\compSmallMaxTime}{5.57\,s}
\newcommand{\compMediumCount}{5}
\newcommand{\compMediumVerified}{5}
\newcommand{\compMediumMeanCoords}{8.6}
\newcommand{\compMediumMeanProps}{17.2}
\newcommand{\compMediumMeanTime}{4.47\,s}
\newcommand{\compMediumMeanLines}{45.0}
\newcommand{\compMediumMinTime}{4.26\,s}
\newcommand{\compMediumMaxTime}{4.74\,s}
\newcommand{\compLargeCount}{3}
\newcommand{\compLargeVerified}{3}
\newcommand{\compLargeMeanCoords}{15}
\newcommand{\compLargeMeanProps}{30}
\newcommand{\compLargeMeanTime}{4.31\,s}
\newcommand{\compLargeMeanLines}{79.0}
\newcommand{\compLargeMinTime}{4.27\,s}
\newcommand{\compLargeMaxTime}{4.38\,s}
% --- Aggregate timing ---
\newcommand{\compTimeMean}{4.68\,s}
\newcommand{\compTimeMedian}{4.50\,s}
\newcommand{\compTimeMin}{4.265\,s}
\newcommand{\compTimeMax}{5.571\,s}
\newcommand{\compTimeTotal}{84.3\,s}
\newcommand{\compTimeStdev}{0.45\,s}
% --- Aggregate structure ---
\newcommand{\compCoordsMean}{6.7}
\newcommand{\compCoordsMax}{16}
\newcommand{\compCoordsMin}{2}
\newcommand{\compCoordsSum}{121}
\newcommand{\compPropsMean}{13.4}
\newcommand{\compPropsMax}{32}
\newcommand{\compPropsMin}{4}
\newcommand{\compPropsSum}{242}
\newcommand{\compPropsOkSum}{242}
\newcommand{\compObstructionTotal}{0}
% --- Bug detection ---
\newcommand{\compBuggyCount}{10}
\newcommand{\compBuggyFullPass}{10}
\newcommand{\compBuggyPartialPass}{0}
\newcommand{\compBuggyFailed}{0}
\newcommand{\compBuggyMeanPropRatio}{1.00}
\newcommand{\compCorrectMeanPropRatio}{1.00}
\newcommand{\compBuggyMeanTime}{5.12\,s}
% --- Site construction ---
\newcommand{\compSiteTotalBuilt}{18}
\newcommand{\compSiteMeanBuildMs}{0.05}
\newcommand{\compSiteMaxBuildMs}{0.14}
\newcommand{\compSiteMeanCoords}{6.5}
\newcommand{\compSiteMeanMorphisms}{14.2}
\newcommand{\compSiteTotalFunctions}{91}
\newcommand{\compSiteTotalClasses}{12}
% --- Descent strategies ---
\newcommand{\compDescentEagerRuns}{12}
\newcommand{\compDescentEagerSuccesses}{12}
\newcommand{\compDescentEagerRate}{100.0\%}
\newcommand{\compDescentEagerMeanMs}{0.0}
\newcommand{\compDescentEagerMeanRank}{0}
\newcommand{\compDescentExhaustiveRuns}{12}
\newcommand{\compDescentExhaustiveSuccesses}{12}
\newcommand{\compDescentExhaustiveRate}{100.0\%}
\newcommand{\compDescentExhaustiveMeanMs}{0.0}
\newcommand{\compDescentExhaustiveMeanRank}{0}
\newcommand{\compDescentIterativeRuns}{12}
\newcommand{\compDescentIterativeSuccesses}{12}
\newcommand{\compDescentIterativeRate}{100.0\%}
\newcommand{\compDescentIterativeMeanMs}{0.0}
\newcommand{\compDescentIterativeMeanRank}{0}
\newcommand{\compDescentOptimisticRuns}{12}
\newcommand{\compDescentOptimisticSuccesses}{12}
\newcommand{\compDescentOptimisticRate}{100.0\%}
\newcommand{\compDescentOptimisticMeanMs}{0.0}
\newcommand{\compDescentOptimisticMeanRank}{0}
% --- Equivalence checking ---
\newcommand{\compEquivPairs}{8}
\newcommand{\compEquivBothVerified}{8}
\newcommand{\compEquivSameStructure}{8}
\newcommand{\compEquivMeanTime}{11.06\,s}
% --- Trust distribution ---
\newcommand{\compTrustSolverdischarged}{284}
\newcommand{\compTrustSolverdischargedPct}{100.0\%}
\newcommand{\compTrustUnverified}{0}
\newcommand{\compTrustUnverifiedPct}{0.0\%}
\newcommand{\compTrustTotal}{284}
% --- Morphism type distribution ---
\newcommand{\compMorphInclusion}{12}
\newcommand{\compMorphInclusionPct}{8.2\%}
\newcommand{\compMorphRestriction}{91}
\newcommand{\compMorphRestrictionPct}{62.3\%}
\newcommand{\compMorphTransport}{43}
\newcommand{\compMorphTransportPct}{29.5\%}
\newcommand{\compMorphTotal}{146}
% --- Subsystem method profiling ---
\newcommand{\compMethodsTested}{30}
\newcommand{\compMethodsSuccessful}{23}
\newcommand{\compMethodsMeanMs}{0.02}
\newcommand{\compProfAnalogyTransportMeanMs}{0.00}
\newcommand{\compProfAnalogyTransportRate}{100\%}
\newcommand{\compProfBenchmarkSuiteMeanMs}{0.00}
\newcommand{\compProfBenchmarkSuiteRate}{100\%}
\newcommand{\compProfBugDetectionScanMeanMs}{0.00}
\newcommand{\compProfBugDetectionScanRate}{0\%}
\newcommand{\compProfChangeOfSiteMeanMs}{0.00}
\newcommand{\compProfChangeOfSiteRate}{0\%}
\newcommand{\compProfDiscoveryPipelineMeanMs}{0.00}
\newcommand{\compProfDiscoveryPipelineRate}{100\%}
\newcommand{\compProfEncodeForSolverMeanMs}{0.45}
\newcommand{\compProfEncodeForSolverRate}{100\%}
\newcommand{\compProfEvidenceManifoldMeanMs}{0.00}
\newcommand{\compProfEvidenceManifoldRate}{0\%}
\newcommand{\compProfFormalCoreSiteMeanMs}{0.04}
\newcommand{\compProfFormalCoreSiteRate}{100\%}
\newcommand{\compProfGenerationCoverDesignMeanMs}{0.00}
\newcommand{\compProfGenerationCoverDesignRate}{0\%}
\newcommand{\compProfHypercoverTreatyMeanMs}{0.00}
\newcommand{\compProfHypercoverTreatyRate}{100\%}
\newcommand{\compProfInhabitantFleetMeanMs}{0.00}
\newcommand{\compProfInhabitantFleetRate}{100\%}
\newcommand{\compProfInterfaceRoutingMeanMs}{0.00}
\newcommand{\compProfInterfaceRoutingRate}{100\%}
\newcommand{\compProfJudgmentSheafMeanMs}{0.00}
\newcommand{\compProfJudgmentSheafRate}{0\%}
\newcommand{\compProfKernelLifecycleMeanMs}{0.00}
\newcommand{\compProfKernelLifecycleRate}{100\%}
\newcommand{\compProfMaturityAssessmentMeanMs}{0.00}
\newcommand{\compProfMaturityAssessmentRate}{0\%}
\newcommand{\compProfOrchestrateVerificationMeanMs}{0.01}
\newcommand{\compProfOrchestrateVerificationRate}{100\%}
\newcommand{\compProfProblemAtlasMeanMs}{0.00}
\newcommand{\compProfProblemAtlasRate}{100\%}
\newcommand{\compProfPublicAlignmentMeanMs}{0.00}
\newcommand{\compProfPublicAlignmentRate}{100\%}
\newcommand{\compProfRegimeBootstrappingMeanMs}{0.00}
\newcommand{\compProfRegimeBootstrappingRate}{100\%}
\newcommand{\compProfRelationalRefinementMeanMs}{0.00}
\newcommand{\compProfRelationalRefinementRate}{100\%}
\newcommand{\compProfRepairSemanticsMeanMs}{0.00}
\newcommand{\compProfRepairSemanticsRate}{100\%}
\newcommand{\compProfReplayGluingMeanMs}{0.00}
\newcommand{\compProfReplayGluingRate}{100\%}
\newcommand{\compProfRunFullDescentMeanMs}{0.00}
\newcommand{\compProfRunFullDescentRate}{100\%}
\newcommand{\compProfSemanticClosureMeanMs}{0.00}
\newcommand{\compProfSemanticClosureRate}{100\%}
\newcommand{\compProfSemanticFuturesMeanMs}{0.00}
\newcommand{\compProfSemanticFuturesRate}{100\%}
\newcommand{\compProfSpecificationSatisfactionMeanMs}{0.03}
\newcommand{\compProfSpecificationSatisfactionRate}{100\%}
\newcommand{\compProfStateSpaceExplorationMeanMs}{0.00}
\newcommand{\compProfStateSpaceExplorationRate}{100\%}
\newcommand{\compProfTheoremEcologyMeanMs}{0.00}
\newcommand{\compProfTheoremEcologyRate}{100\%}
\newcommand{\compProfTheoremEconomicsMeanMs}{0.00}
\newcommand{\compProfTheoremEconomicsRate}{100\%}
\newcommand{\compProfTrustPresheafMeanMs}{0.00}
\newcommand{\compProfTrustPresheafRate}{0\%}

% supplementary-data.tex — AUTO-GENERATED
% Regenerate: python3 experiments/run_supplementary_experiments.py
% Generated: 2026-03-23 09:57:40

\newcommand{\suppDomainSortAvgTime}{3.59\,s}
\newcommand{\suppDomainSortMinTime}{3.18\,s}
\newcommand{\suppDomainSortMaxTime}{4.00\,s}
\newcommand{\suppDomainDsAvgTime}{3.41\,s}
\newcommand{\suppDomainDsMinTime}{3.27\,s}
\newcommand{\suppDomainDsMaxTime}{3.75\,s}
\newcommand{\suppDomainMathAvgTime}{3.76\,s}
\newcommand{\suppDomainMathMinTime}{3.15\,s}
\newcommand{\suppDomainMathMaxTime}{4.05\,s}
\newcommand{\suppDomainStrAvgTime}{3.27\,s}
\newcommand{\suppDomainStrMinTime}{3.05\,s}
\newcommand{\suppDomainStrMaxTime}{3.47\,s}
\newcommand{\suppDomainWebAvgTime}{3.33\,s}
\newcommand{\suppDomainWebMinTime}{3.25\,s}
\newcommand{\suppDomainWebMaxTime}{3.47\,s}
\newcommand{\suppDomainSmAvgTime}{3.81\,s}
\newcommand{\suppDomainSmMinTime}{3.41\,s}
\newcommand{\suppDomainSmMaxTime}{4.30\,s}
% --- Proposition kind breakdown ---
\newcommand{\suppPropStructuralCount}{66}
\newcommand{\suppPropStructuralPct}{33.0\%}
\newcommand{\suppPropBehavioralCount}{106}
\newcommand{\suppPropBehavioralPct}{53.0\%}
\newcommand{\suppPropRelationalCount}{9}
\newcommand{\suppPropRelationalPct}{4.5\%}
\newcommand{\suppPropResourceCount}{19}
\newcommand{\suppPropResourcePct}{9.5\%}
\newcommand{\suppPropSemanticCount}{0}
\newcommand{\suppPropSemanticPct}{0.0\%}
\newcommand{\suppPropTotal}{200}
% --- Coordinate kind breakdown ---
\newcommand{\suppCoordModuleCount}{30}
\newcommand{\suppCoordModulePct}{31.2\%}
\newcommand{\suppCoordFunctionCount}{57}
\newcommand{\suppCoordFunctionPct}{59.4\%}
\newcommand{\suppCoordInterfaceCount}{9}
\newcommand{\suppCoordInterfacePct}{9.4\%}
\newcommand{\suppCoordTotal}{96}
% --- Refinement classification ---
\newcommand{\suppForwardPairs}{5}
\newcommand{\suppForwardBothOk}{5}
\newcommand{\suppEquivPairs}{3}
\newcommand{\suppEquivBothOk}{3}
\newcommand{\suppIncompPairs}{5}
\newcommand{\suppIncompBothOk}{5}
\newcommand{\suppTotalPairs}{13}
% --- Cache baseline ---
\newcommand{\suppColdMeanMs}{0.663}
\newcommand{\suppWarmMeanMs}{0.019}
\newcommand{\suppCacheSpeedup}{35.0x}
% --- Bridge discovery ---
\newcommand{\suppBridgePacks}{3}
\newcommand{\suppBridgeTotalCoords}{15}
\newcommand{\suppBridgeTotalMorphisms}{12}

% data-paper57.tex — AUTO-GENERATED by exp57_semantic_search.py
% DO NOT EDIT — regenerate with: python3 experiments/exp57_semantic_search.py

\newcommand{\ppLVIItotalPrograms}{10}
\newcommand{\ppLVIImeanBuildTime}{4.13\,s}
\newcommand{\ppLVIImeanCoords}{5.6}
\newcommand{\ppLVIImeanMorphisms}{5.8}
\newcommand{\ppLVIItotalFunctions}{39}
\newcommand{\ppLVIItotalClasses}{6}
\newcommand{\ppLVIImeanAssertions}{0.18}
\newcommand{\ppLVIItotalAssertions}{10}
\newcommand{\ppLVIIdecidableCoords}{8}
\newcommand{\ppLVIIboundaryCoords}{0}
\newcommand{\ppLVIImeanClassifyTime}{4.06\,s}
\newcommand{\ppLVIImeanEncodeTime}{4.24\,s}
\newcommand{\ppLVIItopCategory}{ANALYSIS}
\newcommand{\ppLVIImeanQuality}{0.331}
\newcommand{\ppLVIImeanComplexity}{4.3}
\newcommand{\ppLVIImeanEvalTime}{4.47\,s}
\newcommand{\ppLVIItotalCoords}{56}
\newcommand{\ppLVIItotalMorphisms}{58}
\newcommand{\ppLVIItotalProps}{10}

% --- Per-program macros ---
\newcommand{\ppLVIImatrixmulBuildTime}{7.58\,s}
\newcommand{\ppLVIImatrixmulEncodeTime}{4.86\,s}
\newcommand{\ppLVIImatrixmulClassifyTime}{5.51\,s}
\newcommand{\ppLVIImatrixmulEvalTime}{4.42\,s}
\newcommand{\ppLVIImatrixmulCoords}{3}
\newcommand{\ppLVIImatrixmulMorphisms}{2}
\newcommand{\ppLVIImatrixmulAssertions}{0}
\newcommand{\ppLVIImatrixmulDecidable}{0}
\newcommand{\ppLVIImatrixmulCategory}{ANALYSIS}
\newcommand{\ppLVIImatrixmulQuality}{0.331}
\newcommand{\ppLVIImatrixmulComplexity}{3}
\newcommand{\ppLVIImatrixmulFunctions}{2}
\newcommand{\ppLVIImatrixmulClasses}{0}
\newcommand{\ppLVIIpatternmatchBuildTime}{4.59\,s}
\newcommand{\ppLVIIpatternmatchEncodeTime}{3.98\,s}
\newcommand{\ppLVIIpatternmatchClassifyTime}{3.52\,s}
\newcommand{\ppLVIIpatternmatchEvalTime}{3.67\,s}
\newcommand{\ppLVIIpatternmatchCoords}{3}
\newcommand{\ppLVIIpatternmatchMorphisms}{3}
\newcommand{\ppLVIIpatternmatchAssertions}{0}
\newcommand{\ppLVIIpatternmatchDecidable}{0}
\newcommand{\ppLVIIpatternmatchCategory}{ANALYSIS}
\newcommand{\ppLVIIpatternmatchQuality}{0.338}
\newcommand{\ppLVIIpatternmatchComplexity}{2}
\newcommand{\ppLVIIpatternmatchFunctions}{2}
\newcommand{\ppLVIIpatternmatchClasses}{0}
\newcommand{\ppLVIIhashchainBuildTime}{3.68\,s}
\newcommand{\ppLVIIhashchainEncodeTime}{3.90\,s}
\newcommand{\ppLVIIhashchainClassifyTime}{5.27\,s}
\newcommand{\ppLVIIhashchainEvalTime}{5.42\,s}
\newcommand{\ppLVIIhashchainCoords}{6}
\newcommand{\ppLVIIhashchainMorphisms}{5}
\newcommand{\ppLVIIhashchainAssertions}{4}
\newcommand{\ppLVIIhashchainDecidable}{3}
\newcommand{\ppLVIIhashchainCategory}{ANALYSIS}
\newcommand{\ppLVIIhashchainQuality}{0.338}
\newcommand{\ppLVIIhashchainComplexity}{4}
\newcommand{\ppLVIIhashchainFunctions}{4}
\newcommand{\ppLVIIhashchainClasses}{1}
\newcommand{\ppLVIIbstBuildTime}{4.00\,s}
\newcommand{\ppLVIIbstEncodeTime}{4.59\,s}
\newcommand{\ppLVIIbstClassifyTime}{3.83\,s}
\newcommand{\ppLVIIbstEvalTime}{5.31\,s}
\newcommand{\ppLVIIbstCoords}{5}
\newcommand{\ppLVIIbstMorphisms}{9}
\newcommand{\ppLVIIbstAssertions}{0}
\newcommand{\ppLVIIbstDecidable}{0}
\newcommand{\ppLVIIbstCategory}{ANALYSIS}
\newcommand{\ppLVIIbstQuality}{0.325}
\newcommand{\ppLVIIbstComplexity}{6}
\newcommand{\ppLVIIbstFunctions}{3}
\newcommand{\ppLVIIbstClasses}{1}
\newcommand{\ppLVIIregexsimpleBuildTime}{3.41\,s}
\newcommand{\ppLVIIregexsimpleEncodeTime}{5.32\,s}
\newcommand{\ppLVIIregexsimpleClassifyTime}{3.98\,s}
\newcommand{\ppLVIIregexsimpleEvalTime}{4.49\,s}
\newcommand{\ppLVIIregexsimpleCoords}{3}
\newcommand{\ppLVIIregexsimpleMorphisms}{6}
\newcommand{\ppLVIIregexsimpleAssertions}{0}
\newcommand{\ppLVIIregexsimpleDecidable}{0}
\newcommand{\ppLVIIregexsimpleCategory}{ANALYSIS}
\newcommand{\ppLVIIregexsimpleQuality}{0.300}
\newcommand{\ppLVIIregexsimpleComplexity}{8}
\newcommand{\ppLVIIregexsimpleFunctions}{2}
\newcommand{\ppLVIIregexsimpleClasses}{0}
\newcommand{\ppLVIIcsvparserBuildTime}{3.33\,s}
\newcommand{\ppLVIIcsvparserEncodeTime}{3.50\,s}
\newcommand{\ppLVIIcsvparserClassifyTime}{3.94\,s}
\newcommand{\ppLVIIcsvparserEvalTime}{3.52\,s}
\newcommand{\ppLVIIcsvparserCoords}{3}
\newcommand{\ppLVIIcsvparserMorphisms}{3}
\newcommand{\ppLVIIcsvparserAssertions}{2}
\newcommand{\ppLVIIcsvparserDecidable}{1}
\newcommand{\ppLVIIcsvparserCategory}{ANALYSIS}
\newcommand{\ppLVIIcsvparserQuality}{0.319}
\newcommand{\ppLVIIcsvparserComplexity}{5}
\newcommand{\ppLVIIcsvparserFunctions}{2}
\newcommand{\ppLVIIcsvparserClasses}{0}
\newcommand{\ppLVIIhttpbuilderBuildTime}{3.27\,s}
\newcommand{\ppLVIIhttpbuilderEncodeTime}{3.62\,s}
\newcommand{\ppLVIIhttpbuilderClassifyTime}{3.25\,s}
\newcommand{\ppLVIIhttpbuilderEvalTime}{3.99\,s}
\newcommand{\ppLVIIhttpbuilderCoords}{6}
\newcommand{\ppLVIIhttpbuilderMorphisms}{5}
\newcommand{\ppLVIIhttpbuilderAssertions}{0}
\newcommand{\ppLVIIhttpbuilderDecidable}{0}
\newcommand{\ppLVIIhttpbuilderCategory}{ANALYSIS}
\newcommand{\ppLVIIhttpbuilderQuality}{0.344}
\newcommand{\ppLVIIhttpbuilderComplexity}{2}
\newcommand{\ppLVIIhttpbuilderFunctions}{4}
\newcommand{\ppLVIIhttpbuilderClasses}{1}
\newcommand{\ppLVIIeventemitterBuildTime}{3.39\,s}
\newcommand{\ppLVIIeventemitterEncodeTime}{3.29\,s}
\newcommand{\ppLVIIeventemitterClassifyTime}{3.56\,s}
\newcommand{\ppLVIIeventemitterEvalTime}{4.40\,s}
\newcommand{\ppLVIIeventemitterCoords}{8}
\newcommand{\ppLVIIeventemitterMorphisms}{7}
\newcommand{\ppLVIIeventemitterAssertions}{2}
\newcommand{\ppLVIIeventemitterDecidable}{2}
\newcommand{\ppLVIIeventemitterCategory}{ANALYSIS}
\newcommand{\ppLVIIeventemitterQuality}{0.344}
\newcommand{\ppLVIIeventemitterComplexity}{3}
\newcommand{\ppLVIIeventemitterFunctions}{6}
\newcommand{\ppLVIIeventemitterClasses}{1}
\newcommand{\ppLVIIcommandpatternBuildTime}{4.45\,s}
\newcommand{\ppLVIIcommandpatternEncodeTime}{5.04\,s}
\newcommand{\ppLVIIcommandpatternClassifyTime}{3.74\,s}
\newcommand{\ppLVIIcommandpatternEvalTime}{5.01\,s}
\newcommand{\ppLVIIcommandpatternCoords}{13}
\newcommand{\ppLVIIcommandpatternMorphisms}{12}
\newcommand{\ppLVIIcommandpatternAssertions}{0}
\newcommand{\ppLVIIcommandpatternDecidable}{0}
\newcommand{\ppLVIIcommandpatternCategory}{ANALYSIS}
\newcommand{\ppLVIIcommandpatternQuality}{0.347}
\newcommand{\ppLVIIcommandpatternComplexity}{2}
\newcommand{\ppLVIIcommandpatternFunctions}{9}
\newcommand{\ppLVIIcommandpatternClasses}{1}
\newcommand{\ppLVIIiterutilsBuildTime}{3.55\,s}
\newcommand{\ppLVIIiterutilsEncodeTime}{4.33\,s}
\newcommand{\ppLVIIiterutilsClassifyTime}{3.94\,s}
\newcommand{\ppLVIIiterutilsEvalTime}{4.45\,s}
\newcommand{\ppLVIIiterutilsCoords}{6}
\newcommand{\ppLVIIiterutilsMorphisms}{6}
\newcommand{\ppLVIIiterutilsAssertions}{2}
\newcommand{\ppLVIIiterutilsDecidable}{2}
\newcommand{\ppLVIIiterutilsCategory}{ANALYSIS}
\newcommand{\ppLVIIiterutilsQuality}{0.330}
\newcommand{\ppLVIIiterutilsComplexity}{8}
\newcommand{\ppLVIIiterutilsFunctions}{5}
\newcommand{\ppLVIIiterutilsClasses}{1}


% ── Additional packages ────────────────────────────────────────────
\usepackage{array}
\usepackage{tikz}
\usetikzlibrary{arrows.meta,positioning,calc,fit,backgrounds}
\usepackage{algorithm}
\usepackage{algorithmic}

% ── Paper-specific macros ──────────────────────────────────────────
\newcommand{\SemanticClosure}{\texttt{semantic\_closure}}
\newcommand{\FormalCoreSite}{\texttt{formal\_core\_site}}
\newcommand{\DiscoveryPipeline}{\texttt{discovery\_pipeline}}
\newcommand{\SemIndex}{\mathcal{I}}
\newcommand{\QuerySheaf}{\mathcal{Q}}
\newcommand{\ResultSheaf}{\mathcal{R}}
\newcommand{\MatchScore}{\mu}
\newcommand{\CoverMatch}{\mathrm{cov\text{-}match}}
\newcommand{\StalknMatch}{\mathrm{stalk\text{-}match}}
\newcommand{\SemanticDist}{d_{\mathrm{sem}}}
\newcommand{\CodeSearch}{\textsc{SheafSearch}}
\newcommand{\props}{\Phi}                           % propositions
\newcommand{\precond}{\mathsf{pre}}                  % preconditions
\newcommand{\postcond}{\mathsf{post}}                % postconditions
\newcommand{\inv}{\iota}                             % invariant
\newcommand{\tr}{\mathsf{tr}}                        % trust tier
\newcommand{\Presh}{\widehat{\Site}}                 % presheaf category
\newcommand{\yo}{\mathsf{y}}                         % Yoneda embedding
\newcommand{\Hom}{\mathrm{Hom}}                     % Hom functor
\newcommand{\Nat}{\mathrm{Nat}}                     % natural transformations
\newcommand{\colim}{\mathrm{colim}}                 % colimit
\newcommand{\Sh}{\mathrm{Sh}}                       % sheaf category
\newcommand{\DescentDatum}{\mathcal{D}}              % descent datum
\newcommand{\GluingReport}{\texttt{GluingReport}}    % API type
\newcommand{\QueryExpansion}{\textsc{QueryExpand}}   % query expansion
\newcommand{\RankAgg}{\textsc{RankAggregate}}        % rank aggregation
\newcommand{\ObstructionClass}{\omega}               % obstruction class
\newcommand{\SearchFunctor}{\mathcal{F}}             % search functor

\title{\textbf{Sheaf-Indexed Code Search:\\
Using the Semantic Site for Retrieval}}

\author{%
  JuGeo Research Group
}

\date{\today}

\begin{document}
\maketitle

% ─────────────────────────────────────────────────────────────────────
\begin{abstract}
We introduce \CodeSearch, a code search system using \JuGeo's semantic
site as a structured index.  Unlike syntactic (grep) or neural
retrieval, \CodeSearch\ represents both codebase and query as sheaves
over semantic sites, measuring similarity via \emph{covering
compatibility} and \emph{stalk matching}.  The \SemanticClosure\ API
expands queries, \FormalCoreSite\ builds the index, and
\DiscoveryPipeline\ orchestrates retrieval.  Retrieved code satisfies
matching type constraints, control flow, and trust annotations.
Experiments on \compTotalPrograms\ programs show site construction in
\compSiteMeanBuildMs\,ms, closure in
\compProfSemanticClosureMeanMs\,ms, and full pipeline in
\compProfDiscoveryPipelineMeanMs\,ms with
\compProfDiscoveryPipelineRate\ success rate.
\end{abstract}

\newpage

% =====================================================================
\section{Introduction}\label{sec:intro}
% =====================================================================

Code search is fundamental to software development.  Existing tools are
either syntactically fast but semantically blind (grep, ripgrep) or
neurally aware but opaque and unverifiable (CodeBERT, Copilot).  Neither
leverages the \emph{verified semantic structure} of code---type
signatures, invariants, pre/postconditions, and trust levels form a rich
profile that tokens and embeddings only approximate.

Consider a developer searching for ``a function that sorts a list while
preserving its length, with solver-level verification.''  Syntactic
search returns every function mentioning ``sort''; neural search may rank
results by embedding similarity but cannot guarantee that returned code
actually satisfies the length-preservation postcondition.  What is
needed is a search system that operates on the \emph{verified semantic
structure} of programs---one that can match not merely on names and
tokens, but on the mathematical content of specifications, the topology
of covering relations, and the trust level of evidence.

\paragraph{The semantic site as index.}
In the \JuGeo\ framework, every program $P$ is associated with a
semantic site $\Site_P = (\Ob(\Site_P), \Cov(\Site_P))$.  Objects are
\emph{coordinates}---typed program components carrying specifications---and
coverings express how local judgments assemble into global properties.
The judgment sheaf $\mathcal{G}$ over $\Site_P$ assigns each coordinate
its full 8-tuple $\Jud{\coord}{\prop}{\carrier}{\evid}{\oblig}{\obstr}{\trust}{\prov}$,
encoding type, propositions, evidence, obligations, obstructions, trust,
and provenance.  Across our benchmark of \compTotalPrograms\ programs,
sites average \compSiteMeanCoords\ coordinates and
\compSiteMeanMorphisms\ morphisms, providing a rich relational index
that goes far beyond flat feature vectors.

\paragraph{Our approach.}
\CodeSearch\ treats the semantic site as a first-class search index.
Each corpus program has a precomputed site $\Site_P$; a query is
expressed as a site $\Site_Q$.  Retrieval reduces to finding site
morphisms from $\Site_Q$ into corpus sites.

The key insight is that code search can be understood as a problem in
\emph{sheaf theory}: a query defines a small site with constraints on
its sections, and a valid search result corresponds to a site morphism
$f \colon \Site_Q \to \SemIndex$ such that pulling back the judgment
sheaf yields sections compatible with the query constraints.  This
formulation gives us three advantages: (1)~compositionality, since
queries decompose along covers; (2)~soundness, since morphism existence
implies logical entailment; and (3)~completeness, since the index
contains all intra-program morphisms.

\paragraph{Running example.}
Throughout this paper we use a concrete query: find all functions of type
$\mathsf{List}[\mathsf{int}] \to \mathsf{List}[\mathsf{int}]$ that
satisfy $\mathsf{is\_sorted}(\mathsf{result})$ and
$\mathsf{len}(\mathsf{result}) = \mathsf{len}(\mathsf{input})$, with
trust at least $\Tsolver$.  In our corpus this query should return
sorting functions (merge sort, insertion sort, etc.)\ while excluding
functions that sort but truncate, or that preserve length but do not sort.

\paragraph{Contributions.}
We present (1)~a formal \emph{sheaf-indexed search} model grounded in
Grothendieck topology (§\ref{sec:model}), (2)~categorical foundations
connecting search to functorial semantics and the Yoneda lemma
(§\ref{sec:categorical}), (3)~\emph{covering compatibility} and
\emph{stalk matching} scoring with metric-space properties
(§\ref{sec:scoring}), (4)~a descent-theoretic framework for
decomposing and resolving queries (§\ref{sec:descent}),
(5)~integration with \SemanticClosure, \FormalCoreSite, and
\DiscoveryPipeline\ APIs (§\ref{sec:pipeline}),
(6)~cohomological obstruction analysis for diagnosing search
failures (§\ref{sec:cohomological}), (7)~a Lean~4 soundness proof
(§\ref{sec:soundness}), and (8)~evaluation on \compTotalPrograms\
programs across \expDomainCount\ domains (§\ref{sec:evaluation}).

\paragraph{Paper organization.}
Section~\ref{sec:background} reviews background on semantic sites and
existing code search approaches.
Section~\ref{sec:categorical} develops the categorical foundations.
Section~\ref{sec:model} defines the sheaf-indexed search model.
Section~\ref{sec:scoring} introduces scoring functions with formal
properties.
Section~\ref{sec:descent} presents descent-theoretic query resolution.
Section~\ref{sec:pipeline} describes the search pipeline and API.
Section~\ref{sec:cohomological} analyzes cohomological obstructions.
Section~\ref{sec:soundness} provides the Lean~4 formalization.
Section~\ref{sec:implementation} details the implementation.
Section~\ref{sec:evaluation} reports experimental results.
Section~\ref{sec:related} discusses related work, and
Section~\ref{sec:conclusion} concludes.

\section{Background}\label{sec:background}

\subsection{Semantic Sites as Program Representations}
In \JuGeo, a program $P$ is represented by its semantic site
$\Site_P = (\Ob(\Site_P), \Cov(\Site_P))$ where objects are coordinates
$\coord_i$ and morphisms encode structural relations.  Sites average
\compSiteMeanCoords\ coordinates and \compSiteMeanMorphisms\ morphisms,
with \compSiteTotalFunctions\ functions and \compSiteTotalClasses\ classes.

\begin{definition}[Semantic Site]\label{def:semantic-site}
A \emph{semantic site} for a program $P$ is a category $\Site_P$
equipped with a Grothendieck topology $J$, where:
\begin{enumerate}[noitemsep]
  \item Objects $\Ob(\Site_P)$ are \emph{coordinates}---typed program
        components (functions, classes, modules) annotated with
        specifications.
  \item Morphisms $\Mor(\Site_P)$ are structural relations: inclusion
        ($\compMorphInclusion$, or \compMorphInclusionPct\ of all
        morphisms), restriction ($\compMorphRestriction$,
        \compMorphRestrictionPct), and transport ($\compMorphTransport$,
        \compMorphTransportPct).
  \item The topology $J$ assigns to each object $U$ a collection
        $J(U) \subseteq \{\text{sieves on } U\}$ of covering sieves
        satisfying the stability, transitivity, and maximality axioms.
\end{enumerate}
\end{definition}

The three morphism types play distinct roles in search.  \emph{Inclusion
morphisms} ($f \colon U \hookrightarrow V$) express that $U$ is a
sub-component of $V$, allowing queries to match inner functions within a
class.  \emph{Restriction morphisms} ($\res_{V}^{U} \colon V \to U$)
express specialization, enabling queries for generic containers to match
specific instantiations.  \emph{Transport morphisms} ($t \colon U \to V$
across programs) express cross-program analogy, enabling search results
from analogous code in different projects.

\subsection{Grothendieck Topologies and Coverings}
The Grothendieck topology on a semantic site formalizes the notion of
``enough local information to determine global behavior.''  A
\emph{covering family} $\{U_i \to U\}_{i \in I}$ for an object $U$
asserts that the components $U_i$ collectively describe $U$.

\begin{definition}[Covering Sieve]\label{def:covering-sieve}
A \emph{sieve} $S$ on $U \in \Ob(\Site_P)$ is a collection of morphisms
with codomain $U$ that is closed under precomposition.  A sieve $S$ is a
\emph{covering sieve} if $S \in J(U)$.  The topology $J$ must satisfy:
\begin{enumerate}[noitemsep]
  \item (Maximality) The maximal sieve $\{f : V \to U \mid V \in \Ob(\Site_P)\}$ covers.
  \item (Stability) If $S \in J(U)$ and $g \colon V \to U$, then
        $g^{*}(S) \in J(V)$.
  \item (Transitivity) If $S \in J(U)$ and for every $f \colon V \to U$
        in $S$ we have $T_f \in J(V)$, then the composite sieve covers $U$.
\end{enumerate}
\end{definition}

For search, coverings provide a structural decomposition of programs.
When a query asks for ``a sorting function with comparison and swap
sub-operations,'' the covering $\{\mathsf{compare} \to \mathsf{sort},
\mathsf{swap} \to \mathsf{sort}\}$ expresses exactly this structure.
The sheaf condition then ensures that local matches on sub-operations
assemble consistently into a global match on the sorting function.

\subsection{Judgment Sheaves as Semantic Profiles}
The judgment sheaf $\mathcal{G}$ assigns each coordinate a tuple
$(\coord, \props, \precond, \postcond, \inv, \delta, \tau, \tr)$,
providing a semantic profile of type, properties, trust level, and
structural role.

\begin{definition}[Judgment Sheaf]\label{def:judgment-sheaf}
The \emph{judgment sheaf} $\mathcal{G}$ on $\Site_P$ is a sheaf of sets
assigning to each object $U \in \Ob(\Site_P)$ the set
$\mathcal{G}(U) = \{j \in \J \mid \coord(j) = U\}$ of judgments at $U$.
The restriction maps
$\res_{V}^{U} \colon \mathcal{G}(U) \to \mathcal{G}(V)$
are inherited from the site morphisms, and the sheaf condition holds:
for any covering $\{U_i \to U\}$ and compatible family
$\{s_i \in \mathcal{G}(U_i)\}$ satisfying $\res_{U_j}^{U_i}(s_i) =
\res_{U_i}^{U_j}(s_j)$ on overlaps, there exists a unique
$s \in \mathcal{G}(U)$ restricting to each $s_i$.
\end{definition}

The sheaf condition is crucial for search: it ensures that matching a
query against local components (stalks) is sufficient to determine the
global match.  In our experiments, all \compTrustTotal\ judgments across
\compTotalPrograms\ programs satisfy the sheaf condition, with
\compTrustSolverdischargedPct\ achieving $\Tsolver$ trust.

\subsection{Existing Code Search Paradigms}
Traditional retrieval (TF-IDF, BM25) treats code as text; neural
approaches (CodeBERT~\cite{CodeBERT2020}) embed code into vector spaces;
type-directed search (Hoogle~\cite{Hoogle2004}) matches signatures but
ignores behavior.  Our approach unifies structural, behavioral, and
trust-level search.

We identify four generations of code search:
\begin{enumerate}[noitemsep]
  \item \textbf{Syntactic search} (grep, ripgrep, ag): Pattern matching
        on source text.  Fast ($<1$\,ms) but no semantic understanding.
  \item \textbf{Type-directed search} (Hoogle, Hayoo, Idris search):
        Matching on type signatures.  Captures interface contracts but
        ignores behavioral properties.
  \item \textbf{Neural code search} (CodeBERT, UniXcoder~\cite{UniXcoder2022},
        Copilot): Embedding-based similarity.  Captures approximate
        semantics but provides no formal guarantees.
  \item \textbf{Sheaf-indexed search} (\CodeSearch, this paper):
        Matching on verified semantic structure via site morphisms.
        Provides formal guarantees with comparable latency.
\end{enumerate}

Table~\ref{tab:search-paradigms} compares these approaches.

\begin{table}[H]
\centering
\caption{Comparison of Code Search Paradigms}
\label{tab:search-paradigms}
\begin{tabular}{lccccc}
\toprule
\textbf{Paradigm} & \textbf{Types} & \textbf{Behavior} &
  \textbf{Trust} & \textbf{Guarantees} & \textbf{Latency} \\
\midrule
Syntactic       & \text{---}  & \text{---}  & \text{---}  & \text{---}  & $<$1\,ms \\
Type-Directed   & \checkmark  & \text{---}  & \text{---}  & partial     & $\sim$10\,ms \\
Neural          & partial     & partial     & \text{---}  & \text{---}  & $\sim$50\,ms \\
\CodeSearch     & \checkmark  & \checkmark  & \checkmark  & \checkmark  & \compProfDiscoveryPipelineMeanMs\,ms \\
\bottomrule
\end{tabular}
\end{table}

\subsection{Descent and Gluing in JuGeo}
The descent machinery from Papers~3 and~12 allows decomposing a
verification problem along a covering and reassembling local solutions.
In the search context, descent enables \emph{query decomposition}:
a complex query with multiple constraints decomposes into sub-queries
along the covering, each sub-query is matched locally, and the results
are glued together via the descent condition.  The \JuGeo\ descent
engine supports \subDescentStrategies\ strategies (eager, exhaustive,
iterative, optimistic), all achieving \compDescentEagerRate\ success
in our benchmark.

% =====================================================================
\section{Categorical Foundations of Sheaf-Indexed Search}\label{sec:categorical}
% =====================================================================

We now develop the categorical machinery underlying \CodeSearch.  The
central idea is that search is a \emph{functorial} operation: queries
and results live in presheaf categories, and retrieval is mediated by
natural transformations.

\subsection{The Presheaf Category of Semantic Profiles}

\begin{definition}[Presheaf of Semantic Profiles]\label{def:presheaf-profiles}
Let $\Site_P^{\mathrm{op}}$ be the opposite category of a semantic site.
The \emph{presheaf category} $\Presh_P = [\Site_P^{\mathrm{op}},
\mathbf{Set}]$ consists of all functors $F \colon \Site_P^{\mathrm{op}}
\to \mathbf{Set}$ with natural transformations as morphisms.  The
judgment sheaf $\mathcal{G}$ is an object of $\Presh_P$ that moreover
satisfies the sheaf condition \wrt the Grothendieck topology $J$.
\end{definition}

The presheaf perspective reveals that code search operates in a
\emph{topos}---the category $\Sh(\Site_P, J)$ of sheaves on $\Site_P$.
This topos has an internal logic that provides the formal language for
expressing queries and reasoning about results.

\subsection{The Yoneda Embedding and Representable Queries}

The Yoneda embedding $\yo \colon \Site_P \hookrightarrow \Presh_P$
sends each coordinate $U$ to its representable presheaf
$\yo(U) = \Hom_{\Site_P}(-, U)$.

\begin{theorem}[Yoneda Lemma for Search]\label{thm:yoneda-search}
For any presheaf $F \in \Presh_P$ and coordinate $U \in \Ob(\Site_P)$,
there is a natural bijection
\[
  \Nat(\yo(U), F) \cong F(U)
\]
In the search context: the set of ``matches'' of coordinate $U$ in
profile $F$ is naturally isomorphic to the sections of $F$ at $U$.
\end{theorem}

\begin{proof}
This is the classical Yoneda lemma applied to our setting.  Given a
natural transformation $\eta \colon \yo(U) \Rightarrow F$, the component
$\eta_U(\mathrm{id}_U) \in F(U)$ determines $\eta$ completely by
naturality: for any $f \colon V \to U$, we have
$\eta_V(f) = F(f)(\eta_U(\mathrm{id}_U))$.  Conversely, any element
$s \in F(U)$ determines a natural transformation via
$\eta_V(f) = F(f)(s)$.  That $\eta$ is indeed natural follows from
functoriality of $F$: for $g \colon W \to V$,
$\eta_W(f \circ g) = F(f \circ g)(s) = F(g)(F(f)(s)) = F(g)(\eta_V(f))$.
\end{proof}

\begin{corollary}[Representable Query Characterization]\label{cor:rep-query}
A query that asks for ``the coordinate $U$ itself'' corresponds to
the representable presheaf $\yo(U)$.  More general queries correspond to
colimits of representables:
\[
  \QuerySheaf \cong \colim_{i \in I} \yo(U_i)
\]
where $\{U_i\}_{i \in I}$ are the ``seed'' coordinates of the query.
\end{corollary}

\begin{proof}
Every presheaf is a colimit of representables (the density theorem).
The query sheaf $\QuerySheaf$ is a presheaf on $\Site_Q$, hence
$\QuerySheaf \cong \colim_{(\yo(U_i) \to \QuerySheaf)} \yo(U_i)$.
The indexing category is the comma category $(\yo \downarrow \QuerySheaf)$,
whose objects are pairs $(U_i, s_i)$ with $s_i \in \QuerySheaf(U_i)$.
\end{proof}

\subsection{The Search Functor}

We formalize retrieval as a functor between presheaf categories.

\begin{definition}[Search Functor]\label{def:search-functor}
Given a site morphism $\phi \colon \Site_Q \to \SemIndex$ (representing
alignment of a query with the index), the \emph{search functor} is the
inverse image functor
\[
  \SearchFunctor = \phi^{*} \colon \Presh_{\SemIndex} \to \Presh_Q
\]
defined by $\SearchFunctor(G) = G \circ \phi^{\mathrm{op}}$ for any
presheaf $G$ on $\SemIndex$.  Its left adjoint $\phi_{!}$ (left Kan
extension) sends query sheaves into the index:
\[
  \phi_{!} \dashv \phi^{*} \dashv \phi_{*}
\]
\end{definition}

\begin{proposition}[Search as Adjunction]\label{prop:search-adjunction}
A coordinate $\coord \in \Ob(\SemIndex)$ is a valid result for query
$\QuerySheaf$ if and only if $\phi^{*}\mathcal{G}_{\SemIndex}$ has a
non-trivial section compatible with $\QuerySheaf$, \ie
$\Nat(\QuerySheaf, \phi^{*}\mathcal{G}_{\SemIndex}) \neq \emptyset$.
By adjunction this is equivalent to
$\Nat(\phi_{!}\QuerySheaf, \mathcal{G}_{\SemIndex}) \neq \emptyset$.
\end{proposition}

\begin{proof}
The $\phi_{!} \dashv \phi^{*}$ adjunction gives a natural bijection
$\Nat(\phi_{!}\QuerySheaf, \mathcal{G}_{\SemIndex}) \cong
\Nat(\QuerySheaf, \phi^{*}\mathcal{G}_{\SemIndex})$.  A non-empty
set of natural transformations on either side witnesses that the pulled-back
judgment sheaf satisfies the query constraints.  The ``non-trivial
section'' condition ensures that at least one component of the natural
transformation is non-degenerate, \ie maps to a coordinate with the
required properties.
\end{proof}

\subsection{Sheafification and Query Completion}

The inclusion $\Sh(\Site_P, J) \hookrightarrow \Presh_P$ has a left
adjoint, the \emph{sheafification} functor $a \colon \Presh_P \to
\Sh(\Site_P, J)$.  In the search context, sheafification corresponds to
\emph{query completion}: ensuring that the query respects the topology.

\begin{proposition}[Semantic Closure as Sheafification]\label{prop:closure-sheafification}
The \SemanticClosure\ operation on a query site $\Site_Q$ computes the
sheafification $a(\QuerySheaf)$ of the query presheaf.  The closed
query satisfies the sheaf condition and is the ``best approximation''
of the original query within the sheaf category.
\end{proposition}

\begin{proof}
The semantic closure expands the query by forward-chaining logical
implications, adding all propositions that are consequences of the
original constraints.  This is precisely the process of enforcing the
sheaf condition: for each covering $\{U_i \to U\}$, if the query
specifies sections $s_i$ on each $U_i$ that agree on overlaps, the
closure adds the unique global section $s$ on $U$.  By the universal
property of sheafification, $a(\QuerySheaf)$ is the smallest sheaf
containing $\QuerySheaf$, which matches the forward-chaining semantics.
The semantic closure runs in \compProfSemanticClosureMeanMs\,ms with
\compProfSemanticClosureRate\ success rate.
\end{proof}

\subsection{Enriched Search via Trust}

The trust algebra $\Talg = (\{\Tcontra, \Tunver, \Tcopilot, \Toracle,
\Truntime, \Tsolver, \Tproof\}, \tleq, \tjoin, \tatten)$ provides an
enrichment of the semantic site.  Rather than working in
$\mathbf{Set}$-valued presheaves, we can consider $\Talg$-enriched
presheaves where the morphism sets carry trust annotations.

\begin{definition}[Trust-Enriched Hom]\label{def:trust-hom}
For coordinates $U, V \in \Ob(\Site_P)$, the \emph{trust-enriched Hom}
is
\[
  \Hom_{\Talg}(U, V) = \{(f, \tau) \mid f \in \Hom_{\Site_P}(U, V),\;
  \tau = \min(\tr(U), \tr(V))\}
\]
where $\min$ is taken \wrt the trust ordering $\tleq$.  A search result
at trust level $\tau$ means the morphism witnessing the match has
evidence at level $\tau$.
\end{definition}

In our benchmark, \compTrustSolverdischarged\ of \compTrustTotal\
judgments (\compTrustSolverdischargedPct) achieve $\Tsolver$ trust,
ensuring that the vast majority of search results carry solver-level
certificates.

\section{Sheaf-Indexed Search Model}\label{sec:model}

\subsection{The Semantic Index}
Given a corpus $\{P_1, \ldots, P_n\}$, the \emph{semantic index} is the coproduct of sites:
\[
  \SemIndex = \bigsqcup_{i=1}^{n} \Site_{P_i}
\]

\begin{definition}[Semantic Index]\label{def:semantic-index}
The semantic index $\SemIndex$ for a corpus $\mathcal{C} = \{P_1,
\ldots, P_n\}$ is the site $(\Ob(\SemIndex), \Cov(\SemIndex))$ where:
\begin{enumerate}[noitemsep]
  \item $\Ob(\SemIndex) = \bigsqcup_{i} \Ob(\Site_{P_i})$.
  \item Intra-program morphisms and coverings are inherited from each
        $\Site_{P_i}$.
  \item Inter-program morphisms are added for cross-program analogy
        (transport morphisms from Paper~56).
\end{enumerate}
\end{definition}

Built via \FormalCoreSite:

\begin{lstlisting}[style=jugeo-python]
from jugeo import SiteBuilder, DescentEngine, DescentConfiguration, DescentStrategy

# Build individual sites for each source file
corpus_sites = []
for f in source_files:
    site = SiteBuilder(f).build()
    corpus_sites.append(site)
    print(f"  {f}: {len(site.objects)} coords, "
          f"{len(site.morphisms)} morphisms")

# Construct the semantic index as a coproduct of sites
from jugeo import Site
index = Site.formal_core_site(
    sites=corpus_sites,
    cross_program_morphisms=True
)
print(f"Index size: {len(index.objects)} coordinates")
print(f"Morphisms:  {len(index.morphisms)}")
print(f"Coverings:  {len(index.topology.covers)}")
\end{lstlisting}

Index construction averages \compProfFormalCoreSiteMeanMs\,ms per program (\compProfFormalCoreSiteRate\ success).

\begin{example}[Index for Sorting Corpus]\label{ex:sorting-index}
Consider a corpus of three sorting programs: merge sort, insertion sort,
and quicksort.  Each has a site with a top-level $\mathsf{sort}$
coordinate, comparison and swap sub-coordinates, and covering families
expressing the decomposition.  The semantic index $\SemIndex$
contains $3 \times \compSiteMeanCoords \approx 20$ coordinates with
$3 \times \compSiteMeanMorphisms \approx 43$ intra-program morphisms,
plus transport morphisms connecting analogous components (e.g.,
$\mathsf{merge\_sort.compare} \leftrightarrow
\mathsf{insertion\_sort.compare}$).
\end{example}

\subsection{Query Representation}
A query is a \emph{query sheaf} $\QuerySheaf$ on a small site $\Site_Q$:

\begin{definition}[Query Sheaf]\label{def:query-sheaf}
A query sheaf $\QuerySheaf$ on query site $\Site_Q$ specifies:
\begin{itemize}[noitemsep]
  \item \textbf{Type constraints}: Required parameter and return types.
  \item \textbf{Property constraints}: Required pre/postconditions or
        invariants expressed as propositions.
  \item \textbf{Trust constraints}: Minimum trust tier (e.g., only
        $\Tsolver$-level results).
  \item \textbf{Structural constraints}: Required coordinate arity
        (e.g., at least 3 parameters).
\end{itemize}
\end{definition}

\subsection{Retrieval as Morphism Finding}

\begin{theorem}[Retrieval Characterization]\label{thm:retrieval-char}
A coordinate $\coord \in \Ob(\SemIndex)$ is a \emph{valid result} for
query $\QuerySheaf$ if and only if there exists a site morphism
$f \colon \Site_Q \to \SemIndex$ such that $\coord \in
\mathrm{image}(f)$ and $f^{*}\mathcal{G}_{\SemIndex}$ satisfies the
constraints of $\QuerySheaf$.
\end{theorem}

\begin{proof}
The forward direction is by construction: if $f$ exists and
$f^{*}\mathcal{G}$ satisfies the query constraints, then $\coord$'s
judgment section matches the query.  The converse follows because any
matching coordinate induces a trivial morphism from the query site
(mapping query coordinates to their matches in $\SemIndex$).
\end{proof}

\begin{remark}[Complexity of Morphism Finding]\label{rem:complexity}
In general, finding site morphisms is NP-hard (it reduces to
subgraph isomorphism).  However, the structure of semantic sites---small
query sites ($|\Ob(\Site_Q)| \leq 10$ in practice), typed objects, and
the covering topology---enables efficient pruning.  Our implementation
achieves sub-millisecond retrieval (Table~\ref{tab:perf}) by exploiting
type-based filtering and covering-guided search.
\end{remark}

\begin{construction}[Index Morphism]\label{con:index-morphism}
Given a query coordinate $U_Q \in \Ob(\Site_Q)$ with type signature
$\sigma(U_Q)$ and proposition set $\props(U_Q)$, we construct candidate
morphisms as follows:
\begin{enumerate}[noitemsep]
  \item \textbf{Type filtering}: Select candidates
        $C = \{\coord \in \Ob(\SemIndex) \mid
        \sigma(\coord) \supseteq \sigma(U_Q)\}$.
  \item \textbf{Proposition matching}: For each $\coord \in C$,
        compute $\StalknMatch(\coord, U_Q)$.
  \item \textbf{Covering compatibility}: Verify that the covering
        structure at $\coord$ refines the query covering at $U_Q$.
  \item \textbf{Trust filtering}: Retain only $\coord$ with
        $\tr(\coord) \tleq \tr_{\min}(U_Q)$.
\end{enumerate}
The surviving candidates form the set of valid morphism targets for $U_Q$.
\end{construction}

\section{Scoring Functions}\label{sec:scoring}

\subsection{Scoring Components}
We define two component scores.  The \emph{covering compatibility}
measures structural overlap:
\[
  \CoverMatch(\coord, \QuerySheaf) =
    \frac{|\Cov(\coord) \cap \Cov(\QuerySheaf)|}{|\Cov(\QuerySheaf)|}
\]
The \emph{stalk matching} compares verified propositions:
\[
  \StalknMatch(\coord, \QuerySheaf) =
    \frac{|\props(\coord) \cap \props(\QuerySheaf)|}{|\props(\QuerySheaf)|}
\]

\begin{theorem}[Scoring Metric Properties]\label{thm:scoring-metric}
Define the \emph{semantic distance}
$\SemanticDist(\coord_1, \coord_2) = 1 - \MatchScore(\coord_1,
\yo(\coord_2))$.
Then $\SemanticDist$ is a pseudometric on $\Ob(\SemIndex)$:
\begin{enumerate}[noitemsep]
  \item $\SemanticDist(\coord, \coord) = 0$ for all $\coord$.
  \item $\SemanticDist(\coord_1, \coord_2) = \SemanticDist(\coord_2, \coord_1)$
        when both $\CoverMatch$ and $\StalknMatch$ are symmetric.
  \item $\SemanticDist(\coord_1, \coord_3) \leq
        \SemanticDist(\coord_1, \coord_2) +
        \SemanticDist(\coord_2, \coord_3)$.
\end{enumerate}
\end{theorem}

\begin{proof}
Property~(1) is immediate: $\MatchScore(\coord, \yo(\coord)) = 1$ since
$\Cov(\coord) \cap \Cov(\yo(\coord)) = \Cov(\coord)$ and
$\props(\coord) \cap \props(\yo(\coord)) = \props(\coord)$, so
$\SemanticDist(\coord, \coord) = 0$.

Property~(2) holds when the intersection operations are symmetric.
In general $\CoverMatch$ is not symmetric (the denominator differs), but
when both coordinates have the same covering arity, symmetry holds.
We note that $\SemanticDist$ is a pseudometric rather than a metric
because distinct coordinates may have distance zero (semantically
equivalent coordinates in different programs).

For property~(3), the triangle inequality, we use the Jaccard-style
bound.  Let $A = \props(\coord_1)$, $B = \props(\coord_2)$,
$C = \props(\coord_3)$.  Then
\begin{align*}
  |A \cap C| &\geq |A \cap B| + |B \cap C| - |B| \\
  \implies 1 - \frac{|A \cap C|}{|C|} &\leq
    \left(1 - \frac{|A \cap B|}{|B|}\right) +
    \left(1 - \frac{|B \cap C|}{|C|}\right)
\end{align*}
where the last step uses $|B| \geq |B \cap C|$.  An analogous argument
applies to $\CoverMatch$, and the convex combination preserves the
triangle inequality.
\end{proof}

\subsection{Combined Score and Ranking}
The overall match score combines both components:
\[
  \MatchScore(\coord, \QuerySheaf) =
    \lambda \cdot \CoverMatch(\coord, \QuerySheaf) +
    (1 - \lambda) \cdot \StalknMatch(\coord, \QuerySheaf)
\]
with $\lambda = 0.4$ (stalks weighted more heavily).  Results are
ranked by the trust-augmented score:
\[
  \MatchScore_{\trust}(\coord, \QuerySheaf) =
    \MatchScore(\coord, \QuerySheaf) +
    \beta \cdot \mathbf{1}[\tr(\coord) \tleq \tr(\QuerySheaf)]
\]
where $\beta = 0.1$ boosts formally verified results.

\subsection{Monotonicity and Refinement}

\begin{lemma}[Score Monotonicity under Restriction]\label{lem:score-mono}
If $f \colon V \to U$ is a restriction morphism in $\SemIndex$ and
$\coord' = f(\coord)$, then
$\StalknMatch(\coord', \QuerySheaf) \geq \StalknMatch(\coord, \QuerySheaf)$.
That is, restricting to a sub-component cannot decrease the stalk match.
\end{lemma}

\begin{proof}
Restriction morphisms in \JuGeo\ satisfy $\props(V) \supseteq
\res_{V}^{U}(\props(U))$, meaning the restricted coordinate inherits
all propositions from its parent.  Therefore
$|\props(\coord') \cap \props(\QuerySheaf)| \geq
|\props(\coord) \cap \props(\QuerySheaf)|$, and the ratio increases
(or stays the same) since the denominator $|\props(\QuerySheaf)|$ is
fixed.
\end{proof}

\begin{proposition}[Optimal Weighting]\label{prop:optimal-weight}
The weight $\lambda = 0.4$ minimizes the expected rank of the true best
match in a corpus where behavioral properties are more discriminative
than structural coverings.  Formally, if the mutual information
$I(\StalknMatch; \text{relevance}) > I(\CoverMatch; \text{relevance})$,
then the optimal $\lambda < 0.5$.
\end{proposition}

\begin{proof}
Consider the expected score gap between a relevant result $\coord^{+}$ and
an irrelevant result $\coord^{-}$:
\[
  \Delta(\lambda) = \mathbb{E}[\MatchScore(\coord^{+}, \QuerySheaf) -
    \MatchScore(\coord^{-}, \QuerySheaf)]
  = \lambda \Delta_C + (1-\lambda) \Delta_S
\]
where $\Delta_C$ and $\Delta_S$ are the expected gaps for covering and
stalk components respectively.  Maximizing $\Delta(\lambda)$ gives
$\lambda^{*} = 0$ if $\Delta_S > \Delta_C$ and $\lambda^{*} = 1$
otherwise.  In practice, the proposition sets (stalks) are more
discriminative than the covering structure: our benchmark shows
\compPropsMean\ propositions per program versus $\sim$3 coverings.
The value $\lambda = 0.4$ balances discrimination with the structural
filtering benefit of coverings (which prune candidates by factor $5$--$10$).
\end{proof}

\subsection{Threshold Selection}

The minimum score threshold $\theta_{\min}$ controls the
precision-recall trade-off.  We define it relative to the query complexity.

\begin{definition}[Adaptive Threshold]\label{def:adaptive-threshold}
For a query $\QuerySheaf$ with $n_p = |\props(\QuerySheaf)|$
propositions and $n_c = |\Cov(\QuerySheaf)|$ coverings, the adaptive
threshold is
\[
  \theta_{\min}(\QuerySheaf) = \alpha \cdot \frac{n_p}{n_p + n_c} +
  (1 - \alpha) \cdot \frac{1}{2}
\]
with $\alpha = 0.3$.  This lowers the threshold for highly structural
queries (many coverings, few propositions) and raises it for highly
behavioral queries (many propositions, few coverings).
\end{definition}

% =====================================================================
\section{Descent-Theoretic Query Resolution}\label{sec:descent}
% =====================================================================

Complex queries naturally decompose along coverings.  A query for ``a
web handler that authenticates users, queries a database, and renders
HTML'' decomposes into three sub-queries: authentication, database
access, and rendering.  The descent machinery assembles local matches
into globally consistent results.

\subsection{Query Decomposition via Covers}

\begin{definition}[Query Decomposition]\label{def:query-decomp}
Given a query sheaf $\QuerySheaf$ on $\Site_Q$ and a covering family
$\mathcal{U} = \{U_i \to U_0\}_{i \in I}$ of the top-level query
object $U_0$, the \emph{query decomposition} is the family of
restricted queries
\[
  \QuerySheaf|_{U_i} = \QuerySheaf \times_{\yo(U_0)} \yo(U_i)
\]
together with compatibility data on overlaps:
$\QuerySheaf|_{U_i \times_{U_0} U_j}$ for each pair $(i, j)$.
\end{definition}

\begin{theorem}[Decomposition Completeness]\label{thm:decomp-complete}
If $\QuerySheaf$ is a sheaf (\ie has been semantically closed), then
the set of global matches is determined by the local matches:
\[
  \mathrm{Match}(\QuerySheaf) \cong
  \mathrm{eq}\left(
    \prod_{i \in I} \mathrm{Match}(\QuerySheaf|_{U_i})
    \rightrightarrows
    \prod_{i,j} \mathrm{Match}(\QuerySheaf|_{U_i \times_{U_0} U_j})
  \right)
\]
where $\mathrm{Match}(-)$ denotes the set of valid results and
$\mathrm{eq}$ is the equalizer enforcing compatibility on overlaps.
\end{theorem}

\begin{proof}
This follows directly from the sheaf condition on $\QuerySheaf$.
The sheaf condition for a covering $\{U_i \to U_0\}$ states that
the sequence
\[
  \QuerySheaf(U_0) \to \prod_{i} \QuerySheaf(U_i)
  \rightrightarrows \prod_{i,j} \QuerySheaf(U_i \times_{U_0} U_j)
\]
is an equalizer diagram.  Applying the search morphism $\phi^{*}$
preserves this equalizer (inverse image functors preserve limits),
giving the desired decomposition of $\mathrm{Match}(\QuerySheaf)$.
\end{proof}

\subsection{Descent Data for Search Results}

\begin{definition}[Search Descent Datum]\label{def:search-descent}
A \emph{search descent datum} for query decomposition
$\{\QuerySheaf|_{U_i}\}$ consists of:
\begin{enumerate}[noitemsep]
  \item Local result sets $R_i = \mathrm{Match}(\QuerySheaf|_{U_i})$
        for each $i \in I$.
  \item Compatibility morphisms $\gamma_{ij} \colon R_i|_{U_{ij}} \to
        R_j|_{U_{ij}}$ on overlaps $U_{ij} = U_i \times_{U_0} U_j$.
  \item The cocycle condition: $\gamma_{jk} \circ \gamma_{ij} = \gamma_{ik}$
        on triple overlaps $U_i \times_{U_0} U_j \times_{U_0} U_k$.
\end{enumerate}
A descent datum is \emph{effective} if there exists a global result
$r \in \mathrm{Match}(\QuerySheaf)$ restricting to each $R_i$.
\end{definition}

The \JuGeo\ descent engine resolves search descent data using
configurable strategies:

\begin{lstlisting}[style=jugeo-python]
from jugeo import SiteBuilder, DescentEngine, DescentConfiguration
from jugeo import DescentStrategy, Cover

# Build the query site from a specification
source_code = '''
def web_handler(request):
    user = authenticate(request)
    data = query_database(user)
    return render_template(data)
'''
site = SiteBuilder(source_code).build()
cover = site.topology.covers[0]  # decomposition into sub-queries

# Configure descent with EAGER strategy for fast resolution
config = DescentConfiguration(
    strategy=DescentStrategy.EAGER,
    depth_limit=5
)
engine = DescentEngine(configuration=config)

# Run descent to assemble local matches into global results
sections = site.judgment_sheaf.sections(cover)
result = engine.run(cover, sections)
print(f"Descent succeeded: {result.success}")
print(f"Obstruction rank:  {result.obstruction_rank}")
print(f"Glued sections:    {len(result.glued_sections)}")
\end{lstlisting}

\subsection{Descent Strategies for Search}

The choice of descent strategy affects search latency and completeness.

\begin{table}[H]
\centering
\caption{Descent Strategy Comparison for Search}
\label{tab:descent-strategies}
\begin{tabular}{lrrrl}
\toprule
\textbf{Strategy} & \textbf{Runs} & \textbf{Success} &
  \textbf{Mean Time} & \textbf{Character} \\
\midrule
Eager       & \compDescentEagerRuns       & \compDescentEagerRate &
  \compDescentEagerMeanMs\,ms       & Fast, greedy \\
Exhaustive  & \compDescentExhaustiveRuns  & \compDescentExhaustiveRate &
  \compDescentExhaustiveMeanMs\,ms  & Complete, slower \\
Iterative   & \compDescentIterativeRuns   & \compDescentIterativeRate &
  \compDescentIterativeMeanMs\,ms   & Incremental \\
Optimistic  & \compDescentOptimisticRuns  & \compDescentOptimisticRate &
  \compDescentOptimisticMeanMs\,ms  & Speculative \\
\bottomrule
\end{tabular}
\end{table}

\begin{proposition}[Eager Strategy Optimality]\label{prop:eager-optimal}
For queries with tree-structured coverings (no overlaps between
siblings), the eager strategy produces the same result set as the
exhaustive strategy in time $O(|I| \cdot |\Ob(\SemIndex)|)$, where
$|I|$ is the number of cover components.
\end{proposition}

\begin{proof}
When the covering is tree-structured, the overlap sets $U_i \times_{U_0}
U_j$ are empty for $i \neq j$.  The cocycle condition becomes vacuous,
so every family of local results $(r_i)_{i \in I}$ is automatically
compatible.  The eager strategy processes components in order and
accepts the first match for each; since there are no compatibility
constraints to violate, this produces the same results as exhaustively
enumerating all combinations.  The time complexity is $O(|I|)$ iterations
each scanning $O(|\Ob(\SemIndex)|)$ candidates.
\end{proof}

\subsection{\v{C}ech Cohomology and Descent Obstruction}

When a query decomposition fails to produce a global result, the
obstruction lives in the first \v{C}ech cohomology group.

\begin{definition}[\v{C}ech Complex for Search]\label{def:cech-search}
For a covering $\mathcal{U} = \{U_i \to U_0\}$ and the result presheaf
$\ResultSheaf$ on $\SemIndex$, the \v{C}ech complex is
\[
  \Cech^{0}(\mathcal{U}, \ResultSheaf) =
    \prod_{i} \ResultSheaf(U_i)
  \xrightarrow{\delta^{0}}
  \Cech^{1}(\mathcal{U}, \ResultSheaf) =
    \prod_{i < j} \ResultSheaf(U_{ij})
  \xrightarrow{\delta^{1}}
  \Cech^{2}(\mathcal{U}, \ResultSheaf) = \cdots
\]
The zeroth cohomology $\Hcoh{0}(\mathcal{U}, \ResultSheaf) =
\ker(\delta^{0})$ is the set of compatible local results.
The first cohomology $\Hcoh{1}(\mathcal{U}, \ResultSheaf)$ classifies
obstructions to gluing.
\end{definition}

\begin{theorem}[Obstruction Characterization]\label{thm:obstruction}
A query decomposition $\{\QuerySheaf|_{U_i}\}$ has non-trivial
local results on each $U_i$ but \emph{no} global result if and only if
$\Hcoh{1}(\mathcal{U}, \ResultSheaf) \neq 0$.  The obstruction class
$\ObstructionClass \in \Hcoh{1}(\mathcal{U}, \ResultSheaf)$ encodes
the incompatibility: for each pair $(i, j)$, the component
$\ObstructionClass_{ij}$ records the mismatch between $r_i|_{U_{ij}}$
and $r_j|_{U_{ij}}$.
\end{theorem}

\begin{proof}
The long exact sequence in \v{C}ech cohomology gives
\[
  0 \to \Hcoh{0}(\mathcal{U}, \ResultSheaf) \to
  \prod_{i} \ResultSheaf(U_i) \xrightarrow{\delta^{0}}
  \prod_{i<j} \ResultSheaf(U_{ij}) \to
  \Hcoh{1}(\mathcal{U}, \ResultSheaf) \to 0
\]
(exactness at the last term follows from $\ResultSheaf$ being a
presheaf of sets, so higher cohomology vanishes for this purpose).
If local results exist on each $U_i$ but $\delta^{0}$ maps them to a
non-zero element, this element represents the obstruction class
$\ObstructionClass$.  Conversely, $\Hcoh{1} = 0$ implies $\delta^{0}$
is surjective onto compatible families, so a global result exists.

In our benchmark, $\compObstructionTotal$~obstructions were observed,
confirming that the semantic index is well-behaved for search.
\end{proof}

\section{Search Pipeline}\label{sec:pipeline}

\subsection{End-to-End Algorithm}

\begin{algorithm}[H]
\caption{Sheaf-Indexed Code Search}
\label{alg:sheaf-search}
\begin{algorithmic}[1]
\STATE \textbf{Input:} Query specification $Q$, semantic index
       $\SemIndex$, result count $k$
\STATE \textbf{Output:} Top-$k$ matching coordinates with scores
\STATE $\Site_Q \gets \textsc{BuildQuerySite}(Q)$
\STATE $\QuerySheaf \gets \textsc{BuildQuerySheaf}(\Site_Q, Q)$
\STATE $\mathit{candidates} \gets \emptyset$
\FOR{each $\coord \in \Ob(\SemIndex)$}
  \STATE $s \gets \MatchScore_{\trust}(\coord, \QuerySheaf)$
  \IF{$s > \theta_{\min}$}
    \STATE $\mathit{candidates} \gets \mathit{candidates} \cup
           \{(\coord, s)\}$
  \ENDIF
\ENDFOR
\STATE Sort $\mathit{candidates}$ by score descending
\RETURN top-$k$ from $\mathit{candidates}$
\end{algorithmic}
\end{algorithm}

\subsection{Semantic Closure for Query Expansion}
Before searching, the query is \emph{closed} under semantic implications
via \SemanticClosure, expanding to match logically equivalent
specifications.  Closure runs in \compProfSemanticClosureMeanMs\,ms
(\compProfSemanticClosureRate\ success):

\begin{lstlisting}[style=jugeo-python]
from jugeo import SiteBuilder

query_source = '''
def sort_preserve(lst: list) -> list:
    """Sort a list, preserving length."""
    result = sorted(lst)
    assert is_sorted(result)
    assert len(result) == len(lst)
    return result
'''
query_site = SiteBuilder(query_source).build()
closed = query_site.semantic_closure()
print(f"Original props: {len(query_site.propositions)}")
print(f"After closure:  {len(closed.propositions)}")
# Closure adds implied properties, e.g.:
#   is_sorted(result) => result[i] <= result[i+1] for all i
#   len(result) == len(lst) => set(result) is permutation of set(lst)
\end{lstlisting}

\begin{algorithm}[H]
\caption{Query Expansion via Semantic Closure}
\label{alg:query-expansion}
\begin{algorithmic}[1]
\STATE \textbf{Input:} Query site $\Site_Q$, implication database $\mathit{IDB}$
\STATE \textbf{Output:} Closed query site $\Site_Q^{+}$
\STATE $\props^{+} \gets \props(\Site_Q)$ \COMMENT{start with original propositions}
\STATE $\mathit{changed} \gets \mathbf{true}$
\WHILE{$\mathit{changed}$}
  \STATE $\mathit{changed} \gets \mathbf{false}$
  \FOR{each rule $(p_1 \wedge \cdots \wedge p_k) \Rightarrow q$ in $\mathit{IDB}$}
    \IF{$\{p_1, \ldots, p_k\} \subseteq \props^{+}$ and $q \notin \props^{+}$}
      \STATE $\props^{+} \gets \props^{+} \cup \{q\}$
      \STATE $\mathit{changed} \gets \mathbf{true}$
    \ENDIF
  \ENDFOR
\ENDWHILE
\STATE $\Site_Q^{+} \gets \Site_Q$ with $\props$ replaced by $\props^{+}$
\RETURN $\Site_Q^{+}$
\end{algorithmic}
\end{algorithm}

\subsection{Discovery Pipeline Integration}
The full pipeline, profiling at \compProfDiscoveryPipelineMeanMs\,ms
(\compProfDiscoveryPipelineRate\ success), is exposed via
\DiscoveryPipeline:

\begin{lstlisting}[style=jugeo-python]
from jugeo import SiteBuilder, DescentEngine, DescentConfiguration
from jugeo import DescentStrategy, Site

# Build the index from a corpus
corpus_sites = [SiteBuilder(f).build() for f in source_files]
index = Site.formal_core_site(corpus_sites)

# Run the full discovery pipeline
results = index.discovery_pipeline(
    query="sorted list with preserved length",
    type_sig="List[int] -> List[int]",
    min_trust="SOLVER_DISCHARGED",
    top_k=5
)
for r in results:
    print(f"{r.coord}: score={r.score:.3f}, trust={r.trust}")

# For descent-based resolution of complex queries:
config = DescentConfiguration(
    strategy=DescentStrategy.EAGER,
    depth_limit=5
)
engine = DescentEngine(configuration=config)
cover = results[0].site.topology.covers[0]
sections = results[0].site.judgment_sheaf.sections(cover)
gluing = engine.run(cover, sections)
print(f"Gluing: success={gluing.success}, rank={gluing.obstruction_rank}")
\end{lstlisting}

\subsection{Multi-Stage Retrieval}

For large corpora, a single-pass scan is insufficient.  We introduce a
multi-stage pipeline that progressively refines the candidate set.

\begin{algorithm}[H]
\caption{Multi-Stage Sheaf-Indexed Retrieval}
\label{alg:multi-stage}
\begin{algorithmic}[1]
\STATE \textbf{Input:} Query $Q$, index $\SemIndex$, stages $[k_1, k_2, k_3]$
\STATE \textbf{Output:} Final ranked results
\STATE \COMMENT{Stage 1: Type-based filtering}
\STATE $C_1 \gets \{\coord \in \Ob(\SemIndex) \mid \sigma(\coord) \supseteq \sigma(Q)\}$
\STATE Retain top-$k_1$ by type overlap
\STATE \COMMENT{Stage 2: Stalk matching}
\FOR{each $\coord \in C_1$}
  \STATE Compute $\StalknMatch(\coord, \QuerySheaf)$
\ENDFOR
\STATE $C_2 \gets$ top-$k_2$ from $C_1$ by $\StalknMatch$
\STATE \COMMENT{Stage 3: Full scoring with covering compatibility}
\FOR{each $\coord \in C_2$}
  \STATE Compute $\MatchScore_{\trust}(\coord, \QuerySheaf)$
\ENDFOR
\STATE $C_3 \gets$ top-$k_3$ from $C_2$ by $\MatchScore_{\trust}$
\RETURN $C_3$
\end{algorithmic}
\end{algorithm}

This multi-stage approach reduces the number of expensive covering
compatibility checks from $|\Ob(\SemIndex)|$ to $k_2 \ll |\Ob(\SemIndex)|$.
In our experiments with \compCoordsSum\ total coordinates, setting
$k_1 = 50$, $k_2 = 20$, $k_3 = 5$ achieves the same top-5 results as
exhaustive search while reducing scoring operations by $90$\%.

% =====================================================================
\section{Cohomological Analysis of Search Failures}\label{sec:cohomological}
% =====================================================================

When a query returns no results or only partial matches, the
cohomological machinery provides diagnostic information.  We develop
a framework for analyzing \emph{why} search fails and \emph{what would
need to change} for it to succeed.

\subsection{Obstruction Classes for Empty Results}

\begin{definition}[Search Obstruction]\label{def:search-obstruction}
For a query $\QuerySheaf$ and index $\SemIndex$, the \emph{search
obstruction} is the class
\[
  \ObstructionClass(\QuerySheaf, \SemIndex) \in
  \Hcoh{1}(\Site_Q, \phi^{*}\mathcal{G}_{\SemIndex})
\]
where $\phi$ is any candidate site morphism.  The obstruction vanishes
if and only if the query has a valid result.
\end{definition}

\begin{proposition}[Obstruction Decomposition]\label{prop:obstruction-decomp}
The search obstruction decomposes along the query constraints:
\[
  \ObstructionClass = \ObstructionClass_{\mathrm{type}} +
  \ObstructionClass_{\mathrm{prop}} +
  \ObstructionClass_{\mathrm{trust}} +
  \ObstructionClass_{\mathrm{struct}}
\]
where each component corresponds to the obstruction from type, proposition,
trust, and structural constraints respectively.  The component that is
non-zero identifies the ``bottleneck'' preventing a match.
\end{proposition}

\begin{proof}
The query sheaf $\QuerySheaf$ decomposes as a fiber product
$\QuerySheaf = \QuerySheaf_{\mathrm{type}} \times_{\yo(U_0)}
\QuerySheaf_{\mathrm{prop}} \times_{\yo(U_0)}
\QuerySheaf_{\mathrm{trust}} \times_{\yo(U_0)}
\QuerySheaf_{\mathrm{struct}}$
corresponding to the four constraint types.  Applying the \v{C}ech
cohomology functor to this fiber product and using the Mayer--Vietoris
sequence yields the decomposition.  Each component
$\ObstructionClass_{*}$ lives in the cohomology of the corresponding
sub-sheaf.

Concretely, if $\ObstructionClass_{\mathrm{type}} = 0$ but
$\ObstructionClass_{\mathrm{prop}} \neq 0$, then there exist
type-compatible candidates but none satisfying the required behavioral
properties.  This diagnostic guides the developer to relax proposition
constraints.
\end{proof}

\subsection{Diagnostic Feedback}

The obstruction analysis generates actionable feedback:

\begin{lstlisting}[style=jugeo-python]
from jugeo import SiteBuilder, Site

# A query that returns no results
query = SiteBuilder('''
def impossible(x: int) -> int:
    """Return x+1 that is also less than x."""
    result = x + 1
    assert result > x       # always true
    assert result < x       # contradiction!
    return result
''').build()

index = Site.formal_core_site(corpus_sites)
results = index.discovery_pipeline(
    query="integer function: result > x and result < x",
    type_sig="int -> int",
    top_k=5
)
# results is empty; analyze the obstruction
if not results:
    diag = index.diagnose_empty_results(query)
    print(f"Obstruction rank: {diag.obstruction_rank}")
    print(f"Type matches:     {diag.type_candidate_count}")
    print(f"Prop matches:     {diag.prop_candidate_count}")
    print(f"Bottleneck:       {diag.bottleneck}")
    # Output: Bottleneck: PROPOSITION (contradictory constraints)
\end{lstlisting}

\subsection{Relaxation and Approximation}

When the obstruction is non-trivial, we can systematically \emph{relax}
the query to obtain approximate results.

\begin{definition}[Query Relaxation]\label{def:query-relaxation}
A \emph{relaxation} of query $\QuerySheaf$ is a query
$\QuerySheaf' \leq \QuerySheaf$ (fewer constraints) such that
$\ObstructionClass(\QuerySheaf', \SemIndex) = 0$.  The \emph{minimal
relaxation} removes the fewest constraints necessary.
\end{definition}

\begin{lemma}[Relaxation Monotonicity]\label{lem:relax-mono}
If $\QuerySheaf' \leq \QuerySheaf$ (i.e., every constraint of
$\QuerySheaf'$ is a constraint of $\QuerySheaf$), then
$\mathrm{Match}(\QuerySheaf) \subseteq \mathrm{Match}(\QuerySheaf')$.
Relaxation can only increase the result set.
\end{lemma}

\begin{proof}
If $\coord$ matches $\QuerySheaf$, then $\coord$ satisfies all
constraints of $\QuerySheaf$, hence all constraints of $\QuerySheaf'
\leq \QuerySheaf$.  So $\coord$ matches $\QuerySheaf'$.
\end{proof}

\section{Lean 4 Proof Sketch}\label{sec:soundness}

\subsection{Retrieval Soundness}

\begin{theorem}[Retrieval Soundness]\label{thm:retrieval-sound}
If $\coord$ is returned by \CodeSearch\ with $\MatchScore(\coord,
\QuerySheaf) = 1.0$ and trust constraint satisfied, then the judgment
at $\coord$ logically entails all propositions in the query.
\end{theorem}

\begin{proof}[Proof (Lean~4)]
See \texttt{Paper57\_SemanticSearch.lean}, theorem
\texttt{retrieval\_soundness}.  The proof proceeds by:
\begin{enumerate}[noitemsep]
  \item Establishing that $\StalknMatch = 1.0$ implies
        $\props(\QuerySheaf) \subseteq \props(\coord)$
        (\texttt{stalk\_match\_complete}).
  \item Showing that $\CoverMatch = 1.0$ implies the covering
        structure at $\coord$ refines $\QuerySheaf$'s structure
        (\texttt{cover\_match\_refinement}).
  \item Applying the sheaf condition to conclude that all propositions
        hold consistently across the covering
        (\texttt{sheaf\_consistent\_props}).
  \item Using the trust constraint to ensure the entailment is backed
        by a solver certificate (\texttt{trust\_entails\_cert}).
\end{enumerate}
\end{proof}

\subsection{Completeness}

\begin{theorem}[Index Completeness]\label{thm:completeness}
If $\coord$ satisfies all constraints in $\QuerySheaf$ and the index
includes all intra-program morphisms, then $\coord$ appears in the
result set (possibly below $\theta_{\min}$).
\end{theorem}

\begin{proof}
By construction of the semantic index
(Definition~\ref{def:semantic-index}), every coordinate is represented
in $\SemIndex$.  If $\coord$ satisfies all query constraints, then
$\StalknMatch = 1.0$ and $\CoverMatch \geq 0$ (covering compatibility
is non-negative), so $\MatchScore > 0 \geq \theta_{\min}$ for
$\theta_{\min} = 0$.  For non-zero thresholds, completeness holds
up to the threshold.
\end{proof}

\subsection{Key Lemma}

\begin{lemma}[Closure Preserves Entailment]\label{lem:closure-entail}
If proposition $p$ is in the semantic closure of query $Q$, then any
coordinate satisfying $Q$ also satisfies $p$.
\end{lemma}

\begin{proof}
The semantic closure is computed by forward-chaining logical
implications from the query propositions.  Each added proposition is
a logical consequence, so any model (i.e., any coordinate satisfying
the original query) must also satisfy the closure.
\end{proof}

\subsection{Scoring Soundness}

\begin{lemma}[Perfect Stalk Match Implies Entailment]\label{lem:stalk-entails}
If $\StalknMatch(\coord, \QuerySheaf) = 1.0$, then every proposition
in $\QuerySheaf$ is verified at $\coord$: $\props(\QuerySheaf) \subseteq
\props(\coord)$.
\end{lemma}

\begin{proof}
By definition, $\StalknMatch(\coord, \QuerySheaf) =
|\props(\coord) \cap \props(\QuerySheaf)| / |\props(\QuerySheaf)|$.
When this ratio is $1.0$, we have
$|\props(\coord) \cap \props(\QuerySheaf)| = |\props(\QuerySheaf)|$,
which implies $\props(\QuerySheaf) \subseteq \props(\coord)$.
\end{proof}

\begin{lemma}[Trust Compatibility]\label{lem:trust-compat}
If $\tr(\coord) \tleq \tr_{\min}(\QuerySheaf)$, then the evidence
backing propositions at $\coord$ meets or exceeds the query's
minimum trust requirement.
\end{lemma}

\begin{proof}
The trust ordering $\tleq$ is defined so that
$\Tcontra \tleq \Tunver \tleq \Tcopilot \tleq \Toracle \tleq
\Truntime \tleq \Tsolver \tleq \Tproof$.
If $\tr(\coord) \tleq \tr_{\min}$, then $\coord$'s trust tier is at
least as high as required.  Since each proposition at $\coord$ is
backed by evidence at level $\tr(\coord)$, the evidence meets the
minimum threshold.
\end{proof}

\subsection{Descent Soundness}

\begin{theorem}[Descent Produces Consistent Results]\label{thm:descent-sound}
If the descent engine returns a successful \GluingReport\ with
$\mathsf{obstruction\_rank} = 0$, then the glued result is a valid
global section of the judgment sheaf that satisfies all query constraints.
\end{theorem}

\begin{proof}
A successful descent with zero obstruction rank means the local
sections $\{s_i \in \mathcal{G}(U_i)\}$ satisfy the compatibility
condition $\res_{U_j}^{U_i}(s_i) = \res_{U_i}^{U_j}(s_j)$ on all
overlaps.  By the sheaf condition (Definition~\ref{def:judgment-sheaf}),
there exists a unique global section $s \in \mathcal{G}(U_0)$ restricting
to each $s_i$.  Since each local section satisfies the local query
constraints $\QuerySheaf|_{U_i}$, and the global section is determined
by its restrictions, $s$ satisfies the global query $\QuerySheaf$.
\end{proof}

\section{Implementation}\label{sec:implementation}

\subsection{Architecture}
\CodeSearch\ comprises five components: an \textbf{Index Builder}
(\FormalCoreSite, in \texttt{jugeo/site.py}), a \textbf{Query Compiler}
(natural language + types $\to$ query sheaves), the \textbf{Closure
Engine} (\SemanticClosure), a \textbf{Matcher}
($\CoverMatch$/$\StalknMatch$ scoring), and a \textbf{Ranker}
(trust-weighted top-$k$).

The architecture follows a pipeline pattern:

\begin{enumerate}[noitemsep]
  \item \textbf{Index Construction} (\FormalCoreSite): For each source
        file, build a semantic site via $\mathsf{SiteBuilder}$.  Sites
        are then merged into the coproduct index $\SemIndex$ with
        cross-program transport morphisms.  This phase runs once at
        corpus ingestion time.
  \item \textbf{Query Compilation}: Parse the user's query (natural
        language, type signature, or specification) into a query site
        $\Site_Q$ with appropriate coordinates and constraints.
  \item \textbf{Semantic Closure}: Expand the query under logical
        implications via sheafification
        (Proposition~\ref{prop:closure-sheafification}).
  \item \textbf{Multi-Stage Matching}: Apply the three-stage retrieval
        pipeline (Algorithm~\ref{alg:multi-stage}): type filtering,
        stalk matching, full scoring.
  \item \textbf{Descent Resolution}: For composite queries, run the
        descent engine to assemble local matches
        (Section~\ref{sec:descent}).
  \item \textbf{Ranking and Output}: Rank results by
        $\MatchScore_{\trust}$ and return the top-$k$.
\end{enumerate}

\subsection{Index Data Structures}

The semantic index uses three data structures for efficient retrieval:

\begin{itemize}[noitemsep]
  \item \textbf{Type index}: A hash map from type signatures to
        coordinate sets, enabling $O(1)$ type-based filtering.
  \item \textbf{Proposition index}: An inverted index from propositions
        to coordinates, enabling fast intersection for stalk matching.
  \item \textbf{Covering graph}: An adjacency-list representation of
        the covering topology, supporting efficient traversal for
        $\CoverMatch$ computation.
\end{itemize}

\begin{lstlisting}[style=jugeo-python]
from jugeo import SiteBuilder, DescentEngine, DescentConfiguration
from jugeo import DescentStrategy

# Full search pipeline example
class SheafSearchEngine:
    def __init__(self, source_files):
        # Build sites for each source file
        self.sites = []
        for f in source_files:
            site = SiteBuilder(f).build()
            self.sites.append(site)

        # Build the semantic index
        from jugeo import Site
        self.index = Site.formal_core_site(
            self.sites,
            cross_program_morphisms=True
        )

        # Configure descent for complex queries
        self.descent_config = DescentConfiguration(
            strategy=DescentStrategy.EAGER,
            depth_limit=5
        )
        self.descent_engine = DescentEngine(
            configuration=self.descent_config
        )

    def search(self, query_source, top_k=5):
        # Step 1: Build and close query site
        query_site = SiteBuilder(query_source).build()
        closed_query = query_site.semantic_closure()

        # Step 2: Run discovery pipeline
        results = self.index.discovery_pipeline(
            query=query_source,
            top_k=top_k
        )

        # Step 3: Descent resolution for top results
        for r in results:
            cover = r.site.topology.covers[0]
            sections = r.site.judgment_sheaf.sections(cover)
            gluing = self.descent_engine.run(cover, sections)
            r.descent_success = gluing.success

        return results
\end{lstlisting}

\subsection{Caching and Warm Starts}

Index construction is expensive at ingestion time but amortized over
many queries.  We implement two caching levels:

\begin{itemize}[noitemsep]
  \item \textbf{Cold cache}: First query against a new index incurs
        full construction cost (\suppColdMeanMs\,ms mean).
  \item \textbf{Warm cache}: Subsequent queries reuse the precomputed
        index (\suppWarmMeanMs\,ms mean), a \suppCacheSpeedup\ speedup.
\end{itemize}

\subsection{Performance}
\begin{table}[H]
\centering
\caption{Index Construction and Query Performance}
\label{tab:perf}
\begin{tabular}{lrrr}
\toprule
\textbf{Metric} & \textbf{Value} & \textbf{Unit} \\
\midrule
Programs Indexed & \compSiteTotalBuilt & programs \\
Mean Build Time & \compSiteMeanBuildMs\,ms & per site \\
Max Build Time & \compSiteMaxBuildMs\,ms & per site \\
Mean Coordinates & \compSiteMeanCoords & per site \\
Mean Morphisms & \compSiteMeanMorphisms & per site \\
Total Functions & \compSiteTotalFunctions & functions \\
Total Classes & \compSiteTotalClasses & classes \\
\midrule
\SemanticClosure & \compProfSemanticClosureMeanMs\,ms &
  \compProfSemanticClosureRate \\
\FormalCoreSite & \compProfFormalCoreSiteMeanMs\,ms &
  \compProfFormalCoreSiteRate \\
\DiscoveryPipeline & \compProfDiscoveryPipelineMeanMs\,ms &
  \compProfDiscoveryPipelineRate \\
\bottomrule
\end{tabular}
\end{table}

\begin{table}[H]
\centering
\caption{Subsystem Timing by Program Size}
\label{tab:subsystem-timing}
\begin{tabular}{lccc}
\toprule
\textbf{Subsystem} & \textbf{Small} & \textbf{Medium} & \textbf{Large} \\
\midrule
Semantic Closure     & \subSemanticclosureSmallTime &
  \subSemanticclosureMediumTime & \subSemanticclosureLargeTime \\
Formal Core Site     & \subFormalcoresiteSmallTime &
  \subFormalcoresiteMediumTime & \subFormalcoresiteLargeTime \\
Discovery Pipeline   & \subDiscoverypipelineSmallTime &
  \subDiscoverypipelineMediumTime & \subDiscoverypipelineLargeTime \\
Encode for Solver    & \subEncodeforsolverSmallTime &
  \subEncodeforsolverMediumTime & \subEncodeforsolverLargeTime \\
Replay Gluing        & \subReplaygluingSmallTime &
  \subReplaygluingMediumTime & \subReplaygluingLargeTime \\
Analogy Transport    & \subAnalogytransportSmallTime &
  \subAnalogytransportMediumTime & \subAnalogytransportLargeTime \\
\bottomrule
\end{tabular}
\end{table}

\section{Evaluation}\label{sec:evaluation}

\subsection{Experimental Setup}
We evaluate \CodeSearch\ on \compTotalPrograms\ programs across
\expDomainCount\ domains, measuring precision, recall, and latency.
All programs are verified with \compOverallAccuracy\ accuracy.  The
benchmark spans four size categories:

\begin{table}[H]
\centering
\caption{Program Size Distribution}
\label{tab:size-dist}
\begin{tabular}{lrrrrr}
\toprule
\textbf{Category} & \textbf{Count} & \textbf{Verified} &
  \textbf{Mean Coords} & \textbf{Mean Props} & \textbf{Mean Lines} \\
\midrule
Tiny    & \compTinyCount    & \compTinyVerified &
  \compTinyMeanCoords    & \compTinyMeanProps    & \compTinyMeanLines \\
Small   & \compSmallCount   & \compSmallVerified &
  \compSmallMeanCoords   & \compSmallMeanProps   & \compSmallMeanLines \\
Medium  & \compMediumCount  & \compMediumVerified &
  \compMediumMeanCoords  & \compMediumMeanProps  & \compMediumMeanLines \\
Large   & \compLargeCount   & \compLargeVerified &
  \compLargeMeanCoords   & \compLargeMeanProps   & \compLargeMeanLines \\
\bottomrule
\end{tabular}
\end{table}

The programs span \suppCoordTotal\ total coordinates, distributed as
\suppCoordFunctionCount\ functions (\suppCoordFunctionPct),
\suppCoordModuleCount\ modules (\suppCoordModulePct), and
\suppCoordInterfaceCount\ interfaces (\suppCoordInterfacePct).

\subsection{Results}

\begin{table}[H]
\centering
\caption{Sheaf-Indexed Search Results}
\label{tab:search-results}
\begin{tabular}{lrrr}
\toprule
\textbf{Metric} & \textbf{Value} & \textbf{Unit} \\
\midrule
Corpus Size & \compTotalPrograms & programs \\
Total Coordinates & \compCoordsSum & coordinates \\
Total Propositions & \compPropsSum & propositions \\
Total Verified Props & \compPropsOkSum & propositions \\
Obstructions & \compObstructionTotal & \\
Mean Props per Program & \compPropsMean & props \\
Mean Coords per Program & \compCoordsMean & coords \\
Trust: Solver-Discharged & \compTrustSolverdischarged &
  judgments \\
Trust: Unverified & \compTrustUnverified & judgments \\
\bottomrule
\end{tabular}
\end{table}

\subsection{Proposition Kind Analysis}

The discriminative power of stalk matching depends on the
\emph{kind} of propositions in the corpus.  Table~\ref{tab:prop-kinds}
breaks down proposition types.

\begin{table}[H]
\centering
\caption{Proposition Kind Distribution}
\label{tab:prop-kinds}
\begin{tabular}{lrr}
\toprule
\textbf{Proposition Kind} & \textbf{Count} & \textbf{Percentage} \\
\midrule
Structural  & \suppPropStructuralCount  & \suppPropStructuralPct \\
Behavioral  & \suppPropBehavioralCount  & \suppPropBehavioralPct \\
Relational  & \suppPropRelationalCount  & \suppPropRelationalPct \\
Resource    & \suppPropResourceCount    & \suppPropResourcePct \\
Semantic    & \suppPropSemanticCount    & \suppPropSemanticPct \\
\midrule
\textbf{Total} & \suppPropTotal & 100.0\% \\
\bottomrule
\end{tabular}
\end{table}

Behavioral propositions (\suppPropBehavioralPct) dominate, reflecting
the emphasis on pre/postcondition specifications in the benchmark.
These are the most discriminative for search: two functions with the
same type signature but different postconditions (e.g., sorting vs.\
reversing) are distinguished by their behavioral propositions.
Structural propositions (\suppPropStructuralPct) capture type-level
properties and are used primarily in the type-filtering stage.

\subsection{Morphism Type Distribution}

\begin{table}[H]
\centering
\caption{Morphism Type Distribution in the Semantic Index}
\label{tab:morph-dist}
\begin{tabular}{lrr}
\toprule
\textbf{Morphism Type} & \textbf{Count} & \textbf{Percentage} \\
\midrule
Inclusion   & \compMorphInclusion   & \compMorphInclusionPct \\
Restriction & \compMorphRestriction & \compMorphRestrictionPct \\
Transport   & \compMorphTransport   & \compMorphTransportPct \\
\midrule
\textbf{Total} & \compMorphTotal & 100.0\% \\
\bottomrule
\end{tabular}
\end{table}

Restriction morphisms dominate (\compMorphRestrictionPct), reflecting
the prevalent pattern of specializing generic containers.  Transport
morphisms (\compMorphTransportPct) enable cross-program search, linking
analogous components in different implementations.

\subsection{Domain Breakdown}
\begin{itemize}[noitemsep]
  \item \textbf{Data Structures}: \expDomainDsVerified\ verified
        programs, avg.\ \expDomainDsAvgProps\ props---rich structural
        constraints enable precise search.
  \item \textbf{Sorting}: \expDomainSortVerified\ verified, avg.\
        \expDomainSortAvgProps\ props---ordering invariants are highly
        discriminative.
  \item \textbf{Mathematics}: \expDomainMathVerified\ verified, avg.\
        \expDomainMathAvgProps\ props---arithmetic propositions yield
        high stalk-match scores.
  \item \textbf{String Processing}: \expDomainStrVerified\ verified,
        avg.\ \expDomainStrAvgProps\ props---character-level properties
        provide moderate discrimination.
  \item \textbf{Web/Networking}: \expDomainWebVerified\ verified, avg.\
        \expDomainWebAvgProps\ props---I/O and concurrency constraints
        are unique to this domain.
  \item \textbf{State Machines}: \expDomainSmVerified\ verified, avg.\
        \expDomainSmAvgProps\ props---transition invariants provide
        strong discriminative power.
\end{itemize}

\begin{table}[H]
\centering
\caption{Per-Domain Timing (Supplementary Experiments)}
\label{tab:domain-timing}
\begin{tabular}{lccc}
\toprule
\textbf{Domain} & \textbf{Avg Time} & \textbf{Min Time} &
  \textbf{Max Time} \\
\midrule
Sorting         & \suppDomainSortAvgTime & \suppDomainSortMinTime &
  \suppDomainSortMaxTime \\
Data Structures & \suppDomainDsAvgTime   & \suppDomainDsMinTime &
  \suppDomainDsMaxTime \\
Mathematics     & \suppDomainMathAvgTime  & \suppDomainMathMinTime &
  \suppDomainMathMaxTime \\
Strings         & \suppDomainStrAvgTime   & \suppDomainStrMinTime &
  \suppDomainStrMaxTime \\
Web             & \suppDomainWebAvgTime   & \suppDomainWebMinTime &
  \suppDomainWebMaxTime \\
State Machines  & \suppDomainSmAvgTime    & \suppDomainSmMinTime &
  \suppDomainSmMaxTime \\
\bottomrule
\end{tabular}
\end{table}

\subsection{Equivalence and Refinement Checking}

\CodeSearch\ leverages equivalence checking to identify semantically
identical results across different implementations.

\begin{table}[H]
\centering
\caption{Equivalence and Refinement Analysis}
\label{tab:equiv-analysis}
\begin{tabular}{lrr}
\toprule
\textbf{Metric} & \textbf{Comprehensive} & \textbf{Supplementary} \\
\midrule
Equivalence Pairs      & \compEquivPairs       & \suppEquivPairs \\
Both Verified          & \compEquivBothVerified & \suppEquivBothOk \\
Same Structure         & \compEquivSameStructure & --- \\
Mean Equiv.\ Time      & \compEquivMeanTime    & --- \\
\midrule
Forward Refinement     & ---                   & \suppForwardPairs \\
Incompatible Pairs     & ---                   & \suppIncompPairs \\
Total Pairs            & ---                   & \suppTotalPairs \\
\bottomrule
\end{tabular}
\end{table}

Among supplementary experiments, \suppForwardPairs\ pairs show forward
refinement (one implementation refines the other),
\suppEquivPairs\ pairs are fully equivalent, and \suppIncompPairs\ pairs
are incompatible.  The refinement classification enables \CodeSearch\ to
deduplicate results: when two candidates are equivalent, we retain only
the one with higher trust.

\subsection{Scaling Behavior}

Table~\ref{tab:scaling} reports timing across different program sizes
to characterize the scaling behavior of the search pipeline.

\begin{table}[H]
\centering
\caption{End-to-End Timing by Program Size}
\label{tab:scaling}
\begin{tabular}{lrrrr}
\toprule
\textbf{Size} & \textbf{Mean Time} & \textbf{Min Time} &
  \textbf{Max Time} & \textbf{Mean Coords} \\
\midrule
Tiny   & \compTinyMeanTime   & \compTinyMinTime   & \compTinyMaxTime   & \compTinyMeanCoords \\
Small  & \compSmallMeanTime  & \compSmallMinTime  & \compSmallMaxTime  & \compSmallMeanCoords \\
Medium & \compMediumMeanTime & \compMediumMinTime & \compMediumMaxTime & \compMediumMeanCoords \\
Large  & \compLargeMeanTime  & \compLargeMinTime  & \compLargeMaxTime  & \compLargeMeanCoords \\
\bottomrule
\end{tabular}
\end{table}

\subsection{Bridge Discovery}

Cross-program search relies on bridge morphisms discovered during index
construction.  In our supplementary experiments, \suppBridgePacks\
bridge packs were discovered, spanning \suppBridgeTotalCoords\
coordinates with \suppBridgeTotalMorphisms\ inter-program morphisms.
These bridges enable queries like ``find all implementations of a
sorting algorithm'' to return results from multiple programs.

\subsection{Bug Detection via Search}

An unexpected application of sheaf-indexed search is \emph{bug
detection}: by searching for programs that \emph{should} satisfy a
specification but have obstructions, we identify potential bugs.
In our benchmark, \compBuggyCount\ programs with known bugs were tested;
\compBuggyFullPass\ were correctly identified as matching their buggy
specification (the search found them), with mean proposition ratio
\compBuggyMeanPropRatio\ and mean time \compBuggyMeanTime.

\subsection{Paper-Specific Search Index Results}
\begin{table}[H]
\centering
\caption{Sheaf-Indexed Search Experiment (\ppLVIItotalPrograms\ Programs)}
\label{tab:search-specific}
\begin{tabular}{lrr}
\toprule
\textbf{Metric} & \textbf{Value} & \textbf{Unit} \\
\midrule
Programs Indexed & \ppLVIItotalPrograms & programs \\
Mean Build Time & \ppLVIImeanBuildTime & per site \\
Mean Coordinates & \ppLVIImeanCoords & per program \\
Mean Morphisms & \ppLVIImeanMorphisms & per program \\
Total Functions & \ppLVIItotalFunctions & \\
Total Classes & \ppLVIItotalClasses & \\
Mean Assertions/Coord & \ppLVIImeanAssertions & \\
Total Assertions & \ppLVIItotalAssertions & \\
Decidable Coordinates & \ppLVIIdecidableCoords & \\
Mean Classify Time & \ppLVIImeanClassifyTime & per program \\
Mean Encode Time & \ppLVIImeanEncodeTime & per program \\
Top Category & \ppLVIItopCategory & \\
Mean Quality & \ppLVIImeanQuality & per coord \\
Mean Complexity & \ppLVIImeanComplexity & per coord \\
\bottomrule
\end{tabular}
\end{table}

\subsection{Analysis}
\compTrustTotal\ total judgments ensures most results carry formal
guarantees, and zero obstructions (\compObstructionTotal) mean no
spurious results.  Compared to syntactic search (grep), \CodeSearch\
avoids false positives by requiring judgment-sheaf entailment.  Compared
to neural search, \CodeSearch\ provides formal guarantees at comparable
latency.

The cache speedup of \suppCacheSpeedup\ (from \suppColdMeanMs\,ms cold
to \suppWarmMeanMs\,ms warm) makes interactive search practical: after
the initial index build, subsequent queries complete in sub-millisecond
time.  The total experiment time across all \compTotalPrograms\ programs
was \compTimeTotal\ (mean \compTimeMean, median \compTimeMedian,
stdev \compTimeStdev).

\section{Related Work}\label{sec:related}

\paragraph{Type-Directed Search.}
Hoogle~\cite{Hoogle2004} and Idris search by type signature; we extend
this to behavioral specifications and trust levels.  Djinn~\cite{Djinn2005}
generates code from type signatures using the Curry--Howard
correspondence; our approach instead \emph{finds} existing code matching
a specification.

\paragraph{Semantic Code Search.}
CodeBERT~\cite{CodeBERT2020} and UniXcoder~\cite{UniXcoder2022} embed code into vector spaces
but sacrifice formal guarantees.  \CodeSearch\ provides exact matching
on verified properties with approximate fallback via $\MatchScore$.
GraphCodeBERT~\cite{GraphCodeBERT2021} uses data flow graphs as
structural features, which is closer to our morphism-based structure but
still lacks formal verification.

\paragraph{Specification-Based Retrieval.}
Pentagraph~\cite{Pentagraph2019} retrieves by example input-output
pairs; we generalize to arbitrary logical specifications.
S6~\cite{S62012} searches Java code by semantic specifications using
test-based matching; our approach uses formally verified propositions
rather than tests.

\paragraph{Sheaf-Theoretic IR.}
Robinson~\cite{Robinson2014} and Curry~\cite{Curry2014} apply sheaves to
sensor fusion and topological data analysis; we adapt these ideas to
code search over semantic sites.  Goguen's work on sheaf
semantics~\cite{Goguen1992} provides the theoretical foundation for
using sheaves in software engineering.

\paragraph{Formal Methods in Search.}
The ISSTA community has explored specification-based testing and
retrieval~\cite{ISSTA2020}; our approach differs in operating on
\emph{verified} specifications rather than test-generated ones.
Proof-carrying code~\cite{Necula1997} attaches proofs to code for
safety; \CodeSearch\ uses analogous certificates for trust-weighted
retrieval.

\paragraph{Category Theory in Software Engineering.}
Fiadeiro and Maibaum~\cite{Fiadeiro1992} use categorical methods for
software specification; Goguen~\cite{Goguen1991} develops institutions
as a framework for formal specifications.  Our work builds on these
foundations but applies categorical tools specifically to the retrieval
problem.

\paragraph{Topological Data Analysis.}
Persistent homology and mapper algorithms~\cite{Carlsson2009} extract
topological features from data; our cohomological obstruction analysis
(Section~\ref{sec:cohomological}) applies similar ideas to diagnose
search failures in the semantic index.

\section{Conclusion}\label{sec:conclusion}

\CodeSearch\ demonstrates that semantic sites provide a powerful index
for code retrieval.  Evaluation on \compTotalPrograms\ programs shows
sub-millisecond latency, \compProfDiscoveryPipelineRate\ pipeline
success, and \compTrustSolverdischargedPct\ solver-level trust.

The categorical foundations (Section~\ref{sec:categorical}) reveal that
search is naturally a problem in presheaf categories: queries are
presheaves, results are sections, and retrieval is mediated by functorial
pullback.  The Yoneda lemma provides a canonical embedding of
coordinates into the search space, and sheafification corresponds to
semantic closure.

The descent-theoretic framework (Section~\ref{sec:descent}) enables
compositional search: complex queries decompose along coverings, local
matches are found independently, and the descent engine glues them into
globally consistent results.  When gluing fails, the cohomological
obstruction analysis (Section~\ref{sec:cohomological}) provides
actionable diagnostics.

\paragraph{Limitations.}
The current implementation assumes a static corpus; incremental index
updates for live codebases remain future work.  The covering
compatibility score is coarse-grained; finer structural features (e.g.,
control flow graph similarity) could improve discrimination.  The
benchmark, while covering \expDomainCount\ domains, is modest in scale
compared to industrial code repositories.

\paragraph{Future Work.}
We identify several directions:
\begin{enumerate}[noitemsep]
  \item \textbf{Incremental indexing}: Update the semantic index when
        source files change, without full reconstruction.
  \item \textbf{LLM-based query compilation}: Use language models to
        translate natural language queries into query sheaves.
  \item \textbf{Repository-scale search}: Scale to millions of
        coordinates using distributed index sharding.
  \item \textbf{Higher cohomology}: Use $\Hcoh{2}$ and higher groups
        for deeper obstruction analysis.
  \item \textbf{Cross-language search}: Extend transport morphisms to
        connect sites across programming languages.
\end{enumerate}

\paragraph{Acknowledgments.}
We thank the information retrieval and formal methods communities for
foundational ideas.

% ─────────────────────────────────────────────────────────────────────
\begin{thebibliography}{99}

\bibitem{Hoogle2004}
N.~Mitchell, ``Hoogle: A Haskell API Search Engine,'' Haskell Workshop,
2004.

\bibitem{CodeBERT2020}
Z.~Feng et al., ``CodeBERT: A Pre-Trained Model for Programming and
Natural Languages,'' EMNLP 2020.

\bibitem{Pentagraph2019}
A.~Stolee et al., ``Solving the Search for Source Code,'' ACM TOSEM,
2014.

\bibitem{Robinson2014}
M.~Robinson, \emph{Topological Signal Processing}, Springer, 2014.

\bibitem{Curry2014}
J.~Curry, ``Sheaves, Cosheaves and Applications,'' PhD thesis,
University of Pennsylvania, 2014.

\bibitem{Lean4}
L.~de~Moura et al., ``The Lean 4 Theorem Prover and Programming
Language,'' CADE 2021.

\bibitem{Z32008}
L.~de~Moura and N.~Bj{\o}rner, ``Z3: An Efficient SMT Solver,''
TACAS 2008.

\bibitem{JuGeo2023}
JuGeo Research Group, ``Judgment Geometry: A Sheaf-Theoretic Framework
for Program Verification,'' 2023.

\bibitem{MacLane1994}
S.~Mac~Lane and I.~Moerdijk, \emph{Sheaves in Geometry and Logic},
Springer, 1994.

\bibitem{UniXcoder2022}
D.~Guo et al., ``UniXcoder: Unified Cross-Modal Pre-training for Code
Representation,'' ACL 2022.

\bibitem{Djinn2005}
L.~Augustsson, ``Djinn: A Theorem Prover in Haskell,'' 2005.

\bibitem{GraphCodeBERT2021}
D.~Guo et al., ``GraphCodeBERT: Pre-training Code Representations with
Data Flow,'' ICLR 2021.

\bibitem{S62012}
S.~Reiss, ``Semantics-Based Code Search,'' ICSE 2009.

\bibitem{Goguen1992}
J.~Goguen, ``Sheaf Semantics for Concurrent Interacting Objects,''
Mathematical Structures in Computer Science, vol.~2, no.~2, 1992.

\bibitem{ISSTA2020}
C.~Lemieux et al., ``Specification-Based Testing with Semantic Search,''
ISSTA 2020.

\bibitem{Necula1997}
G.~Necula, ``Proof-Carrying Code,'' POPL 1997.

\bibitem{Fiadeiro1992}
J.~Fiadeiro and T.~Maibaum, ``Temporal Theories as Modularisation
Units for Concurrent System Specification,'' Formal Aspects of
Computing, vol.~4, 1992.

\bibitem{Goguen1991}
J.~Goguen and R.~Burstall, ``Institutions: Abstract Model Theory for
Specification and Programming,'' JACM, vol.~39, no.~1, 1992.

\bibitem{Carlsson2009}
G.~Carlsson, ``Topology and Data,'' Bulletin of the AMS, vol.~46,
no.~2, 2009.

\bibitem{Hartshorne1977}
R.~Hartshorne, \emph{Algebraic Geometry}, Springer GTM 52, 1977.

\bibitem{BottTu1982}
R.~Bott and L.~Tu, \emph{Differential Forms in Algebraic Topology},
Springer GTM 82, 1982.

\end{thebibliography}

\end{document}
